% !TeX program = xelatex
% Standalone source: all evaluation plot coordinates and references are embedded.
\documentclass[11pt,a4paper]{article}
\usepackage{kotex}
\ifdefined\setmainhangulfont
 \IfFontExistsTF{UnBatang.ttf}{
  \setmainhangulfont{UnBatang.ttf}[BoldFont=UnBatangBold.ttf,
   ItalicFont=UnBatang.ttf,BoldItalicFont=UnBatangBold.ttf,
   Script=Hangul,Language=Korean]
 }{}
\fi
\usepackage[margin=24mm,headheight=15pt]{geometry}
\usepackage{amsmath,amssymb,mathtools,bm}
\usepackage{booktabs,tabularx,array,longtable}
\usepackage[dvipsnames]{xcolor}
\usepackage{tikz}
\usetikzlibrary{arrows.meta,positioning,calc}
\usepackage{pgfplots}
\usepgfplotslibrary{groupplots}
\pgfplotsset{compat=1.18}
\usepackage{enumitem,fancyhdr,listings,xurl,float,needspace}
\usepackage[unicode,colorlinks=true,linkcolor=MidnightBlue,
 citecolor=MidnightBlue,urlcolor=MidnightBlue,
 pdftitle={PADS PD/ET 분류를 위한 MultiROCKET-HYDRA 모델 설명서}]{hyperref}
\definecolor{ink}{HTML}{17324D}
\definecolor{softblue}{HTML}{EDF4FA}
\definecolor{softgray}{HTML}{F5F6F8}
\newcommand{\R}{\mathbb{R}}
\newcommand{\ind}{\mathbf{1}}
\newcommand{\PD}{\mathrm{PD}}
\newcommand{\ET}{\mathrm{ET}}
\newcommand{\code}[1]{\texttt{\detokenize{#1}}}
\DeclareMathOperator{\BA}{BA}
\DeclareMathOperator{\MR}{MR}
\DeclareMathOperator{\HY}{HY}
\DeclareMathOperator{\PPV}{PPV}
\DeclareMathOperator{\MPV}{MPV}
\DeclareMathOperator{\MIPV}{MIPV}
\DeclareMathOperator{\diag}{diag}
\DeclareMathOperator*{\argmax}{arg\,max}
\DeclareMathOperator*{\argmin}{arg\,min}
\setlength{\parindent}{0pt}
\setlength{\parskip}{0.4em}
\setlength{\emergencystretch}{2.5em}
\clubpenalty=10000
\widowpenalty=10000
\postdisplaypenalty=10000
\renewcommand{\contentsname}{목차}
\renewcommand{\abstractname}{초록}
\renewcommand{\figurename}{그림}
\renewcommand{\tablename}{표}
\renewcommand{\refname}{참고문헌}
\renewcommand{\arraystretch}{1.15}
\setlist[itemize]{leftmargin=1.5em,itemsep=0.2em,topsep=0.3em}
\setlist[enumerate]{leftmargin=1.8em,itemsep=0.3em,topsep=0.3em}
\numberwithin{equation}{section}
\pagestyle{fancy}
\fancyhf{}
\fancyhead[L]{\small\color{ink}PADS PD/ET 모델 설명서}
\fancyhead[R]{\small MultiROCKET--HYDRA}
\fancyfoot[C]{\thepage}
\renewcommand{\headrulewidth}{0.3pt}
\lstset{basicstyle=\ttfamily\footnotesize,backgroundcolor=\color{softgray},
 frame=single,rulecolor=\color{gray!30},breaklines=true,columns=fullflexible,
 keepspaces=true,showstringspaces=false,aboveskip=0.7em,belowskip=0.7em}
\newenvironment{keybox}
 {\begin{center}\begin{minipage}{0.94\linewidth}\color{ink}\hrule\vspace{0.6em}\small}
 {\vspace{0.5em}\hrule\end{minipage}\end{center}}
\title{\color{ink}\bfseries PADS 데이터의 PD/ET 분류를 위한\\[0.3em]
 MultiROCKET--HYDRA 모델 설명서\\[0.5em]
 \large 입력 표현, 수학적 구조, 학습 절차 및 외부 평가 결과}
\author{PD/ET Classification Project}
\date{2026년 10월 17일\\[0.3em]\small 구현 기준: \code{sota_pd_et}, aeon 1.6.0}

\begin{document}
\maketitle
\thispagestyle{empty}
\begin{abstract}
이 문서는 PADS v1.0.0의 움직임 신호로 파킨슨병
(Parkinson's disease, PD)과 본태성 떨림(essential tremor, ET)을
구별하기 위해 구현한 MultiROCKET--HYDRA 분류기를 설명한다.
참가자별 신호를 두 종류의 합성곱 특징으로 변환하고,
학습 자료에서 추정한 척도로 정규화한 뒤,
클래스 가중 릿지 분류기를 적합한다.
정규화 강도는 참가자 단위 중첩 교차검증에서 선택한다.
기호의 정의, 각 계산의 순서, 학습되는 요소와 고정 요소를 구분하고,
일반적인 알고리즘 원리와 고정 버전의 실제 계산 규약을 함께 설명한다.
완료된 실행의 참가자별 외부 예측으로 혼동행렬, ROC 및
정밀도--재현율 곡선과 폴드별 지표를 제시하고 해석의 한계를 설명한다.
관측한 성능은 이 코호트와 평가 절차에 관한 것이며
PD/ET 문제의 최고 성능이나 외부 기관의 진단 성능을 확립하지 않는다.
\end{abstract}
\begin{keybox}
\textbf{핵심 원칙.}
분석의 독립 단위는 \textbf{참가자 한 명}이다.
출력은 \textbf{ET를 지지하는 실수 결정 점수}이며 확률이 아니다.
일반 시계열 벤치마크의 근거와 이 PD/ET 코호트에서 관측한 성능은 구분한다.
\end{keybox}
\clearpage
\begingroup
\footnotesize
\setlength{\parskip}{0pt}
\tableofcontents
\endgroup
\clearpage

\section{분류 문제와 기호}
\subsection{참가자 선택과 정답}
현재 선택 대상은 PADS v1.0.0 \cite{pads} 참가자 표의 진단 문자열이
\code{Parkinson's} 또는 \code{Essential Tremor}인 사람이다.
PD 276명과 ET 28명, 총 \(N=304\)명을 사용한다.
새 이진 표적은 다음과 같다.
\begin{equation}
 y_i=\begin{cases}
 0,&i\text{번째 참가자의 공개 진단이 PD일 때},\\
 1,&i\text{번째 참가자의 공개 진단이 ET일 때}.
 \end{cases}\label{eq:labels}
\end{equation}
따라서 양성 클래스는 ET이다.
원래 숫자 라벨 \(2\)는 ET 외의 여러 감별 진단도 포함하므로,
그 값만 사용하여 ET를 고르면 다른 질환이 섞인다.
선택과 라벨 부여는 \code{data.py}의 \code{load_cohort}가 담당한다.

모델 입력은 움직임 신호뿐이다.
ID는 파일 연결과 분할에, 진단명은 정답 생성에만 사용한다.
질병 자유 서술, 약물, 나이, 성별, 설문 답변은 예측 특징에 포함하지 않는다.
PD로 배정되었지만 설명에 ET가 언급된 참가자 8명도 공개 라벨을 유지한다.
그 언급은 감사용 표시로만 남기며 재진단 근거로 사용하지 않는다.

\begin{table}[H]
\centering\small
\caption{주요 기호와 차원.}\label{tab:notation}
\begin{tabularx}{\linewidth}{@{}lX@{}}
\toprule 기호 & 정의\\ \midrule
\(i\in\{1,\ldots,N\}\) & 참가자 인덱스.\\
\(X_i\in\R^{C\times T}\) & 참가자 \(i\)의 입력 행렬. 기본 \(C=132,\ T=976\).\\
\(c\in\{1,\ldots,C\}\) & 과제 구간ㆍ손목ㆍ센서ㆍ축으로 정한 채널.\\
\(t\in\{0,\ldots,T-1\}\) & 채널 내부 표본 인덱스.\\
\(\mathcal I,\ n=|\mathcal I|\) & 하나의 모델에 사용한 학습 참가자 집합과 그 크기.\\
\(\phi_{\mathcal I}\) & 학습 집합에서 적합한 특징 변환과 정규화의 합성.\\
\(z_i\in\R^p\) & 분류기가 받는 결합 특징 \(z_i=\phi_{\mathcal I}(X_i)\).\\
\(\beta\in\R^p,\ \beta_0\in\R\) & 릿지 계수와 절편.\\
\(\alpha>0\) & 릿지 계수 크기에 대한 정규화 강도.\\
\(s_{\mathcal I}(X)\in\R\) & 적합한 모델이 입력 \(X\)에 주는 결정 점수.\\
\bottomrule
\end{tabularx}
\end{table}

\subsection{예측 목표의 범위}
우선적인 목표는 공개 진단 라벨을 새로운 참가자의 움직임으로 예측하는 것이다.
공개 라벨을 잘 예측한다는 사실만으로 질병의 원인, 진단의 절대적 정확성,
실제 진료에서의 효용이 증명되지는 않는다.
이 범위는 이후 모든 성능 해석에 동일하게 적용한다.

\section{입력 행렬의 구성}
\subsection{132개 채널의 의미}
공식 파일 \code{preprocessed/movement/ID_ml.bin}을
little-endian float32 배열로 읽어 \((132,976)\)으로 재배열한다.
전체 입력은 \(X\in\R^{304\times132\times976}\)이며,
132개 채널은 다음 순서의 조합이다.
\begin{equation}
 \underbrace{11}_{\text{과제 구간}}\times
 \underbrace{2}_{\text{왼쪽/오른쪽 손목}}\times
 \underbrace{2}_{\text{가속도/각속도}}\times
 \underbrace{3}_{X,Y,Z\text{ 축}}=132.
\end{equation}
보존된 구간은 아래와 같다.
\begin{center}\small
\begin{tabularx}{0.96\linewidth}{@{}lX@{}}
\toprule 번호 & 공식 구간 이름\\ \midrule
1--4 & \code{Relaxed1}, \code{Relaxed2},
 \code{RelaxedTask1}, \code{RelaxedTask2}\\
5--9 & \code{StretchHold}, \code{HoldWeight}, \code{DrinkGlas},
 \code{CrossArms}, \code{TouchNose}\\
10--11 & \code{Entrainment1}, \code{Entrainment2}\\
\bottomrule
\end{tabularx}
\end{center}
\textbf{132개 채널은 동시에 측정한 132개 물리 센서가 아니다.}
다른 과제의 같은 표본 번호도 행렬의 같은 열에 놓인다.
예를 들어 안정 과제의 \(t=100\)과 자세 유지 과제의 \(t=100\)은
실험 전체에서 같은 절대 시각이 아니다.
이 모델은 과제별 정렬 표현을 다변량 분류기에 넣는다.
서로 다른 과제의 채널 합을 물리적 동시 측정값의 합으로 해석하지 않는다.

\subsection{공식 전처리와 실제 시간}
공식 전처리는 긴 과제를 구간으로 나누고 일부 과제를 제외하며,
가속도에서 추정 중력 추세를 제거한 뒤
각 1024표본 구간의 앞 48개 표본을 자른다.
그 결과가 현재 사용하는 976개 표본이다.
프로젝트는 이 파일을 읽으므로 같은 절단이나 중력 보정을 다시 하지 않는다.
가속도 채널의 단위는 중력가속도 단위인 \(\mathrm g\),
각속도 채널의 단위는 \(\mathrm{rad}/\mathrm s\)이다.
추세 제거는 가속도의 물리 단위를 바꾸지 않는다.

원시 데이터의 명목 표본화율은 100 Hz이지만,
현재 입력은 원시 타임스탬프를 복원하거나 간격이 긴 구간을 다시 보간하지 않는다.
따라서 976개 표본을 항상 정확한 9.76초의 균일한 실제 시간과 동일시하지 않는다.
기존 Fourier 시연은 원시 시각을 이용한 재표본화를 포함하므로
그 스펙트럼과 이 입력은 전처리 정의부터 구분해야 한다.

\subsection{무결성 확인}
로더는 참가자 표, 선택한 참가자 JSON, 움직임 파일과
전처리 순서를 정의하는 두 상위 스크립트의 SHA-256을 확인한다.
ID의 유일성, 표와 JSON의 진단 일치, 파일 크기와 모든 선택 수치의 유한성도 검사한다.
이는 의도한 파일을 의도한 형식으로 읽는지 확인하는 절차이며
코호트의 임상적 대표성이나 진단 자체의 타당성을 증명하는 검사는 아니다.

\section{모델의 전체 구조}
학습 집합 \(\mathcal I\)를 고정하면 다음 연산을 수행한다.
\begin{align}
 r_i&=\MR_{\mathcal I}(X_i)\in\R^{p_{\MR}},&
 h_i&=\HY_{\mathcal I}(X_i)\in\R^{p_{\HY}},\\
 \widetilde r_i&=S^{\MR}_{\mathcal I}(r_i),&
 \widetilde h_i&=S^{\HY}_{\mathcal I}(h_i),\\
 z_i&=\phi_{\mathcal I}(X_i)
 =\begin{bmatrix}\widetilde h_i\\\widetilde r_i\end{bmatrix},&
 s_i&=\beta_0+\beta^\top z_i,\\
 \widehat y_i&=\ind\{s_i>0\}.&&
\end{align}
\(S^{\MR}_{\mathcal I}\)와 \(S^{\HY}_{\mathcal I}\)는 서로 다른 정규화 연산이다.
점수가 정확히 \(0\)이면 PD를 예측한다.
여기서 특징의 척도를 바꾸는 정규화(scaling)와
릿지 계수에 벌점을 주는 정규화(regularization)는 서로 다른 연산이다.

\begin{figure}[htbp]
\centering
\begin{tikzpicture}[font=\small,>=Latex,
 box/.style={draw=ink,rounded corners=2pt,align=center,fill=softblue,
 text width=49mm,minimum height=12mm},
 wide/.style={box,text width=61mm},every path/.style={line width=0.6pt}]
\node[wide] (in) {참가자별 입력\\\(X_i\in\R^{132\times976}\)};
\node[box,below=9mm of in,xshift=-32mm] (mr) {MultiROCKET\\원 신호와 차분의 풀링 특징};
\node[box,below=9mm of in,xshift=32mm] (hy) {HYDRA\\그룹별 커널 경쟁 특징};
\node[box,below=6mm of mr] (smr) {학습 집합의 통계로\\일반 특징 표준화};
\node[box,below=6mm of hy] (shy) {학습 집합의 통계로\\제곱근ㆍ영점 보존 정규화};
\coordinate (mid) at ($(smr.south)!0.5!(shy.south)$);
\node[wide,below=10mm of mid] (join) {특징 결합\\\(z_i=[\widetilde h_i;\widetilde r_i]\)};
\node[wide,below=6mm of join] (ridge) {클래스 가중 릿지\\\(s_i=\beta_0+\beta^\top z_i\)};
\node[wide,below=6mm of ridge] (out) {\(s_i>0:\mathrm{ET}\quad s_i\leq0:\mathrm{PD}\)};
\draw[->] (in)--(mr); \draw[->] (in)--(hy);
\draw[->] (mr)--(smr); \draw[->] (hy)--(shy);
\draw[->] (smr.south)--(join.north); \draw[->] (shy.south)--(join.north);
\draw[->] (join)--(ridge); \draw[->] (ridge)--(out);
\end{tikzpicture}
\caption{적합된 모델의 예측 과정. 자료 의존적 파라미터와 정규화 통계는
해당 학습 집합에서만 결정한다.}\label{fig:pipeline}
\end{figure}

\subsection{고정 요소와 적합 대상}
이 구조는 진단 라벨의 손실을 역전파하여 합성곱 가중치를 학습하는
일반적인 심층 CNN과 적합 방식이 다르다.
MultiROCKET은 정해진 커널 계열을 사용하며 학습 응답에서 임계값을 정한다.
HYDRA는 사전에 정한 시드로 무작위 커널과 채널 조합을 만든다.
그 뒤 두 정규화 통계와 릿지 계수를 적합한다.
라벨을 직접 사용하는 지도 학습은 가중 릿지 적합과
내부 검증을 통한 \(\alpha\) 선택이다.

라벨을 사용하지 않는 변환도 검증 자료로 적합하면 평가 절차가 달라진다.
따라서 내부 학습 폴드마다 특징 변환과 정규화기를 모두 새로 적합한다.

\section{MultiROCKET의 특징 추출}
\subsection{원 신호와 1차 차분}
원 신호와 함께 채널별 차분을 사용한다.
\begin{equation}
 (\Delta X_i)_{c,t}=X_{i,c,t+1}-X_{i,c,t},\qquad t=0,\ldots,T-2.
\end{equation}
이는 실제 시간 간격으로 나누지 않은 표본 차분이다.
원 신호가 가속도ㆍ각속도이고 타임스탬프 간격 문제도 있으므로,
차분을 곧바로 물리적 시간 미분이라고 해석하지 않는다.

\subsection{확장 합성곱}
길이 9의 커널 \(a=(a_0,\ldots,a_8)\), 확장 간격 \(d\in\mathbb N\),
선택한 채널 집합 \(\mathcal C\)를 고정하자.
원 신호 또는 차분의 영점 패딩된 입력을 \(U\)라 하면 응답의 개념식은
\begin{equation}
 v_t=\sum_{c\in\mathcal C}\sum_{\ell=0}^{8}
 a_\ell U_{c,t+(\ell-4)d}
\end{equation}
이다. 범위 밖의 값은 \(0\)이다.
실제 풀링에는 설정에 따라 전체 응답 또는 패딩을 제외한 내부 응답을 사용하므로,
아래 풀링 길이는 별도 기호 \(m\)으로 둔다.

기본 커널은 9개 위치 중 3개에 \(2\), 나머지 6개에 \(-1\)을 놓는
\(\binom93=84\)개 패턴이다.
확장 간격은 유효 범위 \(1+8d\)를 바꿔 서로 다른 길이의 모양을 검사한다.
다변량 채널 선택과 확장 간격의 세부 배치는 고정 라이브러리에 위임한다.
임계값 \(b\)는 학습 참가자의 합성곱 응답에서 분위수로 정하므로,
전체 변환은 학습 자료와 무관한 고정 함수가 아니다.
\code{normalise=False}이므로 입력 채널을 미리 각각 단위 분산으로 바꾸지 않는다.

\subsection{네 가지 풀링 통계}
응답 \(v=(v_0,\ldots,v_{m-1})\), 임계값 \(b\)에 대해
\(\mathcal P=\{t:v_t>b\}\), \(q=|\mathcal P|\)라 두면
\begin{align}
 \PPV(v,b)&=\frac qm,\\
 \MIPV(v,b)&=\begin{cases}q^{-1}\sum_{t\in\mathcal P}t,&q>0,\\-1,&q=0.\end{cases}
\end{align}
PPV는 반응의 비율, MIPV는 반응이 발생한 평균 표본 위치를 나타낸다.
MIPV의 단위는 초가 아니라 표본 인덱스이다.
현재 aeon 1.6.0의 다변량 루틴에서 크기 요약 MPV는
\begin{equation}
 \MPV_{\mathrm{impl}}(v,b)
 =\begin{cases}
 q^{-1}\sum_{t\in\mathcal P}(v_t+b),&q>0,\\
 0,&q=0
 \end{cases}\label{eq:mpv}
\end{equation}
로 계산한다.
\textbf{비교 조건은 \(v_t>b\)이지만 합산은 \(v_t+b\)이다.}
식~\eqref{eq:mpv}는 이름에서 추측한 정의가 아닌 실제 의존 코드의 규약이다.
이를 \(v_t-b\)의 평균으로 바꾸면 현재 구현과 다른 설명이 된다.

LSPV는 반응 구간의 지속성과 간격을 요약한다.
경계 및 동률 규약 때문에 단순한 연속 양수 개수와 항상 같지는 않다.
정확한 계산 순서는 부록~\ref{app:lspv}에 제시한다.
일반적인 방법의 동기는 논문을, 위 세부 규약은 고정 소스를 기준으로 한다
\cite{multirocket,aeonmr}.

\subsection{기본 특징 수}
\code{mr_kernels=6250}은 독립적으로 학습하는 커널 6,250개를 뜻하지 않는다.
84개 기본 패턴을 바탕으로 확장 간격과 임계값 조합에 예산을 배분한다.
한 표현에서 조합 수는
\begin{equation}
 B=84\left\lfloor\frac{6250}{84}\right\rfloor=84\times74=6216
\end{equation}
이며 원 신호ㆍ차분 각각에 네 통계를 적용하므로
\begin{equation}
 p_{\MR}=2\times4\times6216=49728.
\end{equation}
이는 현재 버전의 기본 설정에 대한 값이다.
예산이나 버전이 달라지면 출력 차원을 다시 확인한다.

\section{HYDRA의 커널 경쟁 특징}
\subsection{커널 생성과 그룹 내 경쟁}
성분을 독립인 표준 정규분포
\(a_\ell\sim\mathcal N(0,1)\), \(\ell=0,\ldots,8\)에서 생성한 벡터 \(a\)에 대해
평균을 빼고 절댓값 합으로 나눈다.
\begin{equation}
 \bar a=\frac19\sum_{\ell=0}^{8}a_\ell,\qquad
 \widetilde a_\ell=\frac{a_\ell-\bar a}
 {\sum_{r=0}^{8}|a_r-\bar a|}.
\end{equation}
분모가 \(0\)이 아니면 계수 합은 \(0\), 절댓값 합은 \(1\)이다.
연속 정규분포에서 분모가 \(0\)인 사건의 확률은 \(0\)이며,
유한 정밀도 계산에서는 수치의 유한성 검사도 사용한다.

하나의 표현ㆍ확장 간격ㆍ그룹을 고정하고 \(K=8\)개 응답을 \(v_{k,t}\)라 쓰자.
각 표본에서 승리하는 커널 인덱스는
\begin{equation}
 k^+(t)=\argmax_{1\leq k\leq K}v_{k,t},\qquad
 k^-(t)=\argmin_{1\leq k\leq K}v_{k,t}
\end{equation}
이다. 동률은 사용한 PyTorch 연산의 첫 인덱스 규약을 따른다.
각 커널에 누적하는 값은
\begin{align}
 H^+_k&=\sum_t v_{k,t}\ind\{k=k^+(t)\},\\
 H^-_k&=\sum_t\ind\{k=k^-(t)\}
\end{align}
이다. 첫째는 최대 응답값의 가중 누적, 둘째는 최소 응답 횟수의 누적이다.
따라서 둘 다 단순한 승리 횟수라고 설명하면 정확하지 않다.
최소 응답도 반대 부호 패턴과 강하게 일치한 값일 수 있으므로
가장 나쁜 일치라고 단정하지 않는다 \cite{hydra,aeonhydra}.

현재 \(C=132\)에서는 그룹별로 채널 인덱스 8개를 복원 추출하고,
선택된 채널 합에 커널을 적용한다. 채널이 중복될 수도 있다.
이 혼합은 입력 표현 위의 통계적 연산이며
가속도와 각속도를 같은 물리 단위로 합친 계측량이라는 주장은 아니다.

\subsection{그룹 배분과 특징 수}
기본 \code{hydra_groups=64}는 전체 그룹 수이다.
현재 구현은 원 신호에 32개, 차분에 32개를 배정한다.
각 표현에 64개씩 배정한다고 해석하면 특징 수가 두 배로 잘못 계산된다.
입력 길이 \(T=976\)의 확장 간격은
\begin{equation}
 \mathcal D_{\HY}
 =\left\{2^j:0\leq j\leq\left\lfloor\log_2\frac{T-1}{8}\right\rfloor\right\}
 =\{1,2,4,8,16,32,64\}.
\end{equation}
그러므로 기본 HYDRA 특징 수와 결합 특징 수는
\begin{align}
 p_{\HY}
 &=\underbrace2_{\text{표현}}\times\underbrace7_{\text{간격}}
 \times\underbrace{32}_{\text{표현당 그룹}}
 \times\underbrace8_{\text{커널}}\times\underbrace2_{\text{최대/최소}}
 =7168,\\
 p&=p_{\MR}+p_{\HY}=49728+7168=56896.
\end{align}
특징 차원이 늘어도 독립 참가자 수가 증가하는 것은 아니다.

\section{두 분기의 정규화}
\subsection{MultiROCKET 특징의 표준화}
학습 참가자 \(i\in\mathcal I\)의 \(j\)번째 특징을 \(r_{ij}\)라 하자.
\begin{equation}
 \mu_j=\frac1n\sum_{i\in\mathcal I}r_{ij},\qquad
 \sigma_j=\sqrt{\frac1n\sum_{i\in\mathcal I}(r_{ij}-\mu_j)^2}
\end{equation}
를 계산하고 일반적인 비상수 좌표에서는
\(\widetilde r_{ij}=(r_{ij}-\mu_j)/\sigma_j\)를 사용한다.
\code{StandardScaler}는 상수 좌표의 나눗셈 척도를 \(1\)로 두며,
수치적으로 거의 상수인 좌표에도 라이브러리의 안정화 규약을 적용한다.
검증ㆍ테스트에는 학습 평균과 척도를 그대로 사용한다.

한 특징을 100배로 바꾸면 같은 예측을 만드는 계수는 100분의 1이 되고
그 계수의 제곱 벌점도 달라진다.
따라서 특징 척도를 통제하는 것은 릿지 벌점의 의미를 통제하는 데 필요하다.

\subsection{HYDRA의 제곱근 및 영점 보존 척도}
HYDRA 누적 특징 \(h_{ij}\)를
\begin{equation}
 u_{ij}=\sqrt{\max(h_{ij},0)}
\end{equation}
로 바꾼다. 최대 응답 누적은 음수일 수 있으므로,
이 절단은 각 표본 응답이 아닌 누적된 특징 값에 적용한다.
학습 집합에서
\begin{align}
 \mu_j^{\HY}&=\frac1n\sum_{i\in\mathcal I}u_{ij},&
 \rho_j&=\frac1n\sum_{i\in\mathcal I}\ind\{u_{ij}=0\},\\
 s_j&=\sqrt{\frac1{n-1}\sum_{i\in\mathcal I}(u_{ij}-\mu_j^{\HY})^2},&
 \varepsilon_j&=\rho_j^4+10^{-8}
\end{align}
를 구한다. \(n\geq2\)이고 \(s_j\)는 PyTorch 기본값에 따라 분모 \(n-1\)을 쓴다.
최종 변환은
\begin{equation}
 \widetilde h_{ij}
 =\frac{(u_{ij}-\mu_j^{\HY})\ind\{u_{ij}\neq0\}}{s_j+\varepsilon_j}.
\end{equation}
검증ㆍ테스트에도 이 학습 통계를 사용한다.
마스크 때문에 원래 \(u_{ij}=0\)인 좌표는 변환 뒤에도 \(0\)이다.
이는 일반 표준화와 다르며 결과의 평균이 반드시 \(0\),
분산이 반드시 \(1\)이라는 성질을 주장하지 않는다.

\section{클래스 가중 릿지 분류기}
\subsection{클래스 불균형 가중치}
학습 집합에서 클래스 \(c\in\{0,1\}\)의 수를
\(n_c=\sum_{i\in\mathcal I}\ind\{y_i=c\}\)라 둔다.
두 클래스가 모두 존재할 때
\begin{equation}
 w_c=\frac{n}{2n_c}
\end{equation}
를 적용한다. 각 클래스의 총 가중치는
\(n_cw_c=n/2\)로 동일하다.
ET가 적으면 ET 한 명의 손실에 더 큰 가중치가 붙는다.
참가자를 합성하거나 테스트 코호트를 균형화하는 연산은 아니다.
가중치는 모든 릿지 적합의 학습 인원으로 다시 계산한다.

\subsection{목적함수}
\code{RidgeClassifier}의 이진 표현에 맞춰
\(t_i=2y_i-1\)로 둔다. 즉 PD는 \(-1\), ET는 \(+1\)이다.
주어진 \(\alpha>0\)에서
\begin{equation}
 (\widehat\beta_0,\widehat\beta)=
 \argmin_{\beta_0\in\R,\beta\in\R^p}
 \left\{\sum_{i\in\mathcal I}
 w_{y_i}(t_i-\beta_0-z_i^\top\beta)^2+\alpha\|\beta\|_2^2\right\}
 \label{eq:ridge}
\end{equation}
를 푼다. 코드에 맞춰 손실에 \(1/n\)을 붙이지 않았고,
절편 \(\beta_0\)에는 벌점을 적용하지 않는다.
손실을 평균으로 바꾸면 같은 수치의 \(\alpha\)라도 같은 문제가 아닐 수 있다
\cite{skridge}.

새 입력 \(X_\star\)는
\begin{equation}
 s_\star=\widehat\beta_0+\widehat\beta^\top\phi_{\mathcal I}(X_\star),\qquad
 \widehat y_\star=\ind\{s_\star>0\}
\end{equation}
로 예측한다.
계수는 특징 공간에서 선형이지만 \(\phi_{\mathcal I}\)는 임계값, 최대ㆍ최소,
제곱근 등을 포함하므로 입력 신호에 대한 전체 함수는 일반적으로 비선형이다.

\subsection{특징 수가 참가자 수보다 커도 해가 존재하는 이유}
학습 특징의 행렬을 \(Z\in\R^{n\times p}\),
가중치 행렬을 \(W=\diag(w_{y_i})\)라 두자.
가중 평균과 중심화한 행렬ㆍ벡터를 다음과 같이 정의한다.
\begin{align}
 \bar z_w&=\frac{\sum_{i\in\mathcal I}w_{y_i}z_i}
                  {\sum_{i\in\mathcal I}w_{y_i}},&
 \bar t_w&=\frac{\sum_{i\in\mathcal I}w_{y_i}t_i}
                  {\sum_{i\in\mathcal I}w_{y_i}},\\
 (Z_c)_{i,:}&=(z_i-\bar z_w)^\top,&
 (t_c)_i&=t_i-\bar t_w.
\end{align}
행 번호 \(i\)는 이 식에서 학습 참가자의 행에 대응한다.
현재 클래스 가중치에서는 \(\bar t_w=0\)이지만,
\(\bar z_w\)는 일반적으로 \(0\)이 아니므로 절편이 \(0\)이라고 결론 내릴 수는 없다.
아래 \(I_p\)는 \(p\times p\) 단위행렬이고,
\(\|v\|_2^2=\sum_jv_j^2\)는 유클리드 노름의 제곱이다.
그러면 식~\eqref{eq:ridge}의 해는
\begin{equation}
 \widehat\beta=(Z_c^\top WZ_c+\alpha I_p)^{-1}Z_c^\top Wt_c,\qquad
 \widehat\beta_0=\bar t_w-\bar z_w^\top\widehat\beta.
\end{equation}
이는 해의 수학적 표현이며 실제 코드가 역행렬을 직접 계산한다는 뜻은 아니다.
임의의 \(v\neq0\)에 대해
\begin{equation}
 v^\top(Z_c^\top WZ_c+\alpha I_p)v
 =\|W^{1/2}Z_cv\|_2^2+\alpha\|v\|_2^2>0
\end{equation}
이므로 행렬은 양의 정부호이다.
따라서 \(p>n\)이어도 \(\alpha>0\)이면 계수 해는 유일하다.
다만 이 유일성이 새로운 참가자의 예측 성능을 보장하지는 않는다.

\section{참가자 단위 중첩 교차검증}
\subsection{외부 평가와 내부 선택의 분리}
같은 참가자의 손목이나 과제를 다른 폴드에 넣으면 이미 본 사람의 특성을
활용할 수 있다. 따라서 \((132,976)\) 전체를 한 표본으로 유지한다.
PD와 ET를 나눠 층화하고 모든 학습ㆍ검증 집합에 두 클래스가 있는지 확인한다.

전체 인덱스 집합을 \(\mathcal D=\{1,\ldots,N\}\)라 하자.
서로소 외부 테스트 집합 \(\mathcal E_k\) (\(k=1,\ldots,5\))의 합집합은
\(\mathcal D\)이고 외부 학습 집합은
\(\mathcal T_k=\mathcal D\setminus\mathcal E_k\)이다.
\(\mathcal T_k\) 내부에 다시 3폴드를 구성하여
내부 학습 \(\mathcal U_{kj}\), 검증 \(\mathcal V_{kj}\)를 정한다.
\begin{equation}
 \mathcal U_{kj}\cap\mathcal V_{kj}=\varnothing,\qquad
 \mathcal U_{kj}\cup\mathcal V_{kj}=\mathcal T_k,\qquad j=1,2,3.
\end{equation}
\begin{enumerate}
\item 각 \(\mathcal U_{kj}\)에서 두 특징 변환과 정규화기를 새로 적합한다.
\item 그 학습 특징으로 각 \(\alpha\) 후보의 가중 릿지를 적합한다.
\item 적합된 변환을 \(\mathcal V_{kj}\)에 적용해 검증 점수를 얻는다.
\item 세 내부 검증 결과로 \(\alpha\)를 선택한다.
\item \(\mathcal T_k\) 전체에서 특징ㆍ정규화기를 다시 적합하고,
 선택한 \(\alpha\)로 릿지를 적합한다.
\item 마지막에만 \(\mathcal E_k\)를 예측한다.
\end{enumerate}
외부 테스트는 특징 추출 적합, 척도 추정, \(\alpha\) 선택에 쓰지 않는다.

\subsection{선택 기준과 동률 처리}
기본 후보는
\(\mathcal A=\{10^{-3},10^{-2},10^{-1},1,10,10^2,10^3\}\)이다.
내부 검증 균형 정확도를 \(\BA_{\mathcal V_{kj}}(\alpha)\)라 쓰면
\begin{equation}
 A_k(\alpha)=\frac13\sum_{j=1}^{3}\BA_{\mathcal V_{kj}}(\alpha),\qquad
 A_k^{\max}=\max_{\alpha\in\mathcal A}A_k(\alpha)
\end{equation}
를 계산하고
\begin{equation}
 \alpha_k^\star=\max\{\alpha\in\mathcal A:
 |A_k(\alpha)-A_k^{\max}|\leq10^{-12}\}
\end{equation}
를 고른다. 즉 내부 폴드 균형 정확도의 산술평균이 가장 좋은 값을 쓰고,
동률이면 큰 \(\alpha\)를 택한다.
내부 예측을 합친 지표는 진단용 저장값이며 선택 기준으로 쓰지 않는다.
결정 임계값 \(0\)은 내부ㆍ외부 결과를 보고 조정하지 않는다.

\subsection{계산 횟수와 반복}
기본 외부 평가에서 특징 적합은 \(5(3+1)=20\)번,
릿지 적합은 \(5(3\times7+1)=110\)번이다.
한 내부 학습 폴드의 특징은 7개 \(\alpha\)가 공유하지만
서로 다른 학습 폴드는 적합한 변환을 공유하지 않는다.
기본 시드 \code{20261017}에서 반복ㆍ폴드별 파생 시드를 미리 정한다.
가장 좋은 시드를 결과를 보고 고르는 절차는 없다.
반복 수 \(R\)를 늘리면 참가자마다 반복별 외부 예측 하나를 얻지만,
독립 ET 참가자는 여전히 28명이다.

\section{평가 지표와 해석}
\subsection{균형 정확도와 클래스별 재현율}
한 반복에서 참가자별 외부 예측을 하나씩 모은다.
혼동행렬의 행은 실제 클래스, 열은 예측 클래스이며 둘 다 \([\PD,\ET]\) 순서이다.
\begin{equation}
 \begin{array}{c|cc}
 &\widehat y=0\;(\PD)&\widehat y=1\;(\ET)\\ \hline
 y=0\;(\PD)&\mathrm{TN}&\mathrm{FP}\\
 y=1\;(\ET)&\mathrm{FN}&\mathrm{TP}
 \end{array}
\end{equation}
\begin{align}
 \operatorname{Recall}_{\ET}&=\frac{\mathrm{TP}}{\mathrm{TP}+\mathrm{FN}},&
 \operatorname{Recall}_{\PD}&=\frac{\mathrm{TN}}{\mathrm{TN}+\mathrm{FP}},\\
 \BA&=\tfrac12(\operatorname{Recall}_{\ET}+\operatorname{Recall}_{\PD}),&
 \operatorname{Accuracy}&=\frac{\mathrm{TP}+\mathrm{TN}}N.
\end{align}
주 지표는 클래스 두 개를 같은 비중으로 보는 균형 정확도 \(\BA\)이다.
폴드별 지표를 계산할 때는 \(N\)을 해당 폴드의 인원으로 바꾸어 읽는다.
항상 PD라고 예측하는 기준선은
\begin{equation}
 \operatorname{Accuracy}=\frac{276}{304}\approx0.9079,\qquad
 \BA=0.5,\qquad \operatorname{Recall}_{\ET}=0.
\end{equation}
이 수치는 학습 결과가 아닌 클래스 구성에서 계산된다.
높은 일반 정확도 하나만으로 유용한 분류기라고 주장하기 어려운 이유이다.

\subsection{AUROC와 평균 정밀도}
ET 참가자의 점수 \(s_i\), PD 참가자의 점수 \(s_j\)를 비교하면
\begin{equation}
 \operatorname{AUROC}=\frac1{N_{\ET}N_{\PD}}
 \sum_{\substack{i:y_i=1\\j:y_j=0}}
 \left[\ind\{s_i>s_j\}+\tfrac12\ind\{s_i=s_j\}\right].
\end{equation}
여기서 \(N_{\ET}\)와 \(N_{\PD}\)는 지표를 계산하는 집합의 ET 및 PD 인원이다.
따라서 AUROC는 이진 예측으로 바꾸기 전 실수 점수로 계산한다.
임계값을 높은 값부터 내려 재현율이 증가하는 순서로 놓고 \(R_0=0\)으로 두자.
ET 평균 정밀도(AP)는 그 임계값별 정밀도 \(P_\ell\)와 재현율 \(R_\ell\)로
\begin{equation}
 \operatorname{AP}=\sum_\ell(R_\ell-R_{\ell-1})P_\ell
\end{equation}
를 계산하는 계단형 요약이다. 사다리꼴 적분과 항상 같은 값은 아니다.
점수를 전부 같게 주는 비정보적 예측의 AP 기준값은
ET 비율 \(28/304\approx0.0921\)이다.
유한 자료의 무작위 순위가 매번 정확히 이 AP를 준다는 뜻은 아니다.
추가 지표는 다음과 같이 정의한다.
\begin{align}
 \operatorname{Precision}_{\ET}
 &=\frac{\mathrm{TP}}{\mathrm{TP}+\mathrm{FP}},&
 F_{1,\ET}&=\frac{2\mathrm{TP}}{2\mathrm{TP}+\mathrm{FP}+\mathrm{FN}},\\
 F_{1,\PD}&=\frac{2\mathrm{TN}}{2\mathrm{TN}+\mathrm{FP}+\mathrm{FN}},&
 \operatorname{macro}F_1&=\tfrac12(F_{1,\ET}+F_{1,\PD}).
\end{align}
Matthews 상관계수는
\begin{equation}
 \operatorname{MCC}
 =\frac{\mathrm{TP}\,\mathrm{TN}-\mathrm{FP}\,\mathrm{FN}}
 {\sqrt{(\mathrm{TP}+\mathrm{FP})(\mathrm{TP}+\mathrm{FN})
        (\mathrm{TN}+\mathrm{FP})(\mathrm{TN}+\mathrm{FN})}}
\end{equation}
이다. 분모가 \(0\)이면 해당 라이브러리의 정의를 따라 \(0\)으로 처리한다.
ET 예측이 없을 때 ET 정밀도ㆍF1도 \code{zero_division=0} 규약을 따른다.

\subsection{폴드 점수와 불확실성}
폴드별 점수는 서로 다른 모델에서 나온다.
한 반복에서 점수를 합친 AUROC/AP는 모델 간 점수 척도 차이에 영향을 받을 수
있으므로 폴드별 지표와 기술적 평균도 함께 확인한다.
균형 정확도는 고정 임계값에 의한 참가자별 클래스 예측을 합쳐 계산한다.
기본 테스트 폴드에는 ET가 5명 또는 6명뿐이므로
ET 한 명의 예측 변화가 폴드 재현율을 \(1/5\) 또는 \(1/6\)만큼 바꾼다.
학습 집합이 겹치므로 폴드ㆍ반복의 표준편차는 기술적 산포이지 신뢰구간이 아니다.


% BEGIN GENERATED OUTER EVALUATION

\definecolor{evalteal}{HTML}{087F8C}
\definecolor{evalnavy}{HTML}{17324D}
\definecolor{evalgray}{HTML}{97A3B1}
\definecolor{evalorange}{HTML}{B95C16}

\clearpage
\section{완료된 실행의 외부 평가 결과}\label{sec:observed-results}
\subsection{실행 이력과 평가 단위}
이 절은 \code{artifacts/pads_mr_hydra}에 저장된 완료 상태의 실험을 요약한다.
실행 종료 시각은 \code{2026-10-17T00:21:54} (UTC)이며,
외부 5폴드, 내부 3폴드, 반복 1회,
시드 \(20261017\)를 사용하였다.
아래 그림과 표는 반복 \(1\)에서 참가자별로 정확히 한 번 생성한
외부 보류 예측 \(304\)건에 근거한다.
실제 라벨과 \(\widehat y=\ind\{s>0\}\), 외부ㆍ내부 분할의 참가자 배정,
폴드별 예측 파일 및 저장 지표의 일치를 확인하였다.
최종 전체 자료 적합 모델의 학습 예측과 내부 선택 지표는 여기에 포함하지 않는다.

\begin{table}[H]
\centering\small
\caption{반복 1의 통합 외부 예측 지표와 항상 PD 기준선.
F1ㆍAUROCㆍAP는 \(0\)--\(1\) 척도이며, 나머지는 백분율이다.
ET 정밀도와 F1의 분모가 \(0\)인 기준선은 구현 규약에 따라 \(0\)으로 표시한다.}
\label{tab:observed-metrics}
\begin{tabular}{@{}lrr@{}}
\toprule 지표 & MultiROCKET--HYDRA & 항상 PD\\ \midrule
균형 정확도 & 67.13\% & 50.00\%\\
ET 재현율 & 35.71\% & 0.00\%\\
PD 재현율 & 98.55\% & 100.00\%\\
일반 정확도 & 92.76\% & 90.79\%\\
ET 정밀도 & 71.43\% & 0.00\%\\
ET F1 & 0.4762 & 0.0000\\
AUROC & 0.9000 & 0.5000\\
ET 평균 정밀도(AP) & 0.5776 & 0.0921\\

\bottomrule
\end{tabular}
\end{table}

\begin{figure}[H]
\centering
\begin{tikzpicture}[font=\small]
\begin{scope}[xshift=0.0cm]
\node[font=\small\bfseries] at (2.75,0.8) {(a) Number of participants};
\node at (1.85,0) {PD};
\node at (3.6500000000000004,0) {ET};
\node[anchor=east] at (0.8,-1.05) {PD ($n=276$)};
\filldraw[fill=evalteal!95!white,draw=white,line width=1pt] (0.9500000000000001,-1.9500000000000002) rectangle (2.75,-0.15000000000000002);
\node[text=white,font=\large\bfseries] at (1.85,-1.05) {272};
\filldraw[fill=evalteal!9!white,draw=white,line width=1pt] (2.7500000000000004,-1.9500000000000002) rectangle (4.550000000000001,-0.15000000000000002);
\node[text=evalnavy,font=\large\bfseries] at (3.6500000000000004,-1.05) {4};
\node[anchor=east] at (0.8,-2.85) {ET ($n=28$)};
\filldraw[fill=evalteal!13!white,draw=white,line width=1pt] (0.9500000000000001,-3.75) rectangle (2.75,-1.9500000000000002);
\node[text=evalnavy,font=\large\bfseries] at (1.85,-2.85) {18};
\filldraw[fill=evalteal!11!white,draw=white,line width=1pt] (2.7500000000000004,-3.75) rectangle (4.550000000000001,-1.9500000000000002);
\node[text=evalnavy,font=\large\bfseries] at (3.6500000000000004,-2.85) {10};
\node at (2.75,-4.1) {Predicted class};
\node[rotate=90] at (-1.8,-1.95) {True class};
\end{scope}
\begin{scope}[xshift=7.3cm]
\node[font=\small\bfseries] at (2.75,0.8) {(b) Within each true class};
\node at (1.85,0) {PD};
\node at (3.6500000000000004,0) {ET};
\node[anchor=east] at (0.8,-1.05) {PD ($n=276$)};
\filldraw[fill=evalteal!95!white,draw=white,line width=1pt] (0.9500000000000001,-1.9500000000000002) rectangle (2.75,-0.15000000000000002);
\node[text=white,font=\large\bfseries] at (1.85,-1.05) {98.55\%};
\filldraw[fill=evalteal!9!white,draw=white,line width=1pt] (2.7500000000000004,-1.9500000000000002) rectangle (4.550000000000001,-0.15000000000000002);
\node[text=evalnavy,font=\large\bfseries] at (3.6500000000000004,-1.05) {1.45\%};
\node[anchor=east] at (0.8,-2.85) {ET ($n=28$)};
\filldraw[fill=evalteal!64!white,draw=white,line width=1pt] (0.9500000000000001,-3.75) rectangle (2.75,-1.9500000000000002);
\node[text=evalnavy,font=\large\bfseries] at (1.85,-2.85) {64.29\%};
\filldraw[fill=evalteal!39!white,draw=white,line width=1pt] (2.7500000000000004,-3.75) rectangle (4.550000000000001,-1.9500000000000002);
\node[text=evalnavy,font=\large\bfseries] at (3.6500000000000004,-2.85) {35.71\%};
\node at (2.75,-4.1) {Predicted class};
\node[rotate=90] at (-1.8,-1.95) {True class};
\end{scope}
\end{tikzpicture}

\caption{통합 외부 혼동행렬. 행은 실제 진단, 열은 예측 진단이다.
(a)는 참가자 수, (b)는 각 실제 클래스 안의 백분율이다.
고정 임계값 \(s>0\)에서 PD 272/276명과 ET 10/28명을 맞혔다.
색의 진하기는 각 패널 내부 최댓값을 기준으로 하므로 두 패널 사이의 색을
절대 인원 비교에 사용하지 않는다.}\label{fig:evaluation-confusion}
\end{figure}

일반 정확도는 92.76\%이지만 ET 18/28명은 PD로 예측되었다.
따라서 이 실행의 고정 결정 규칙에서는 PD 재현율(98.55\%)과
ET 재현율(35.71\%)이 크게 다르다.
균형 정확도 67.13\%는 항상 PD 기준선보다 17.13 퍼센트포인트 높다.
이는 이 자료ㆍ분할에서 관측한 차이이며 유의성 검정이나 독립 코호트의
개선 보장을 뜻하지 않는다.

\clearpage
\subsection{순위 지표와 고정 임계값의 구분}
\begin{figure}[H]
\centering
\begin{tikzpicture}
\begin{groupplot}[
 group style={group size=2 by 1,horizontal sep=1.35cm},
 width=6.85cm,height=6.05cm,
 xmin=0,xmax=1,ymin=0,ymax=1.03,
 xtick={0,0.2,0.4,0.6,0.8,1},ytick={0,0.2,0.4,0.6,0.8,1},
 axis lines=left,axis line style={evalgray},
 grid=major,grid style={black!8},
 tick label style={font=\footnotesize},
 label style={font=\small},title style={font=\small\bfseries},
 legend style={font=\scriptsize,draw=none,fill=white,fill opacity=0.94,
  text opacity=1,cells={anchor=west}},
 clip=true]
\nextgroupplot[title={(a) ROC},xlabel={False positive rate (PD)},
 ylabel={True positive rate (ET)},
 legend style={at={(0.5,-0.30)},anchor=north}]
\addplot[evalgray,line width=0.6pt] coordinates {(0,0) (0,0.2) (0.01785714286,0.2) (0.01785714286,0.4) (0.01785714286,0.6) (0.03571428571,0.6) (0.05357142857,0.6) (0.05357142857,0.8) (0.07142857143,0.8) (0.07142857143,1) (0.08928571429,1) (0.1071428571,1) (0.125,1) (0.1428571429,1) (0.1607142857,1) (0.1785714286,1) (0.1964285714,1) (0.2142857143,1) (0.2321428571,1) (0.25,1) (0.2678571429,1) (0.2857142857,1) (0.3035714286,1) (0.3214285714,1) (0.3392857143,1) (0.3571428571,1) (0.375,1) (0.3928571429,1) (0.4107142857,1) (0.4285714286,1) (0.4464285714,1) (0.4642857143,1) (0.4821428571,1) (0.5,1) (0.5178571429,1) (0.5357142857,1) (0.5535714286,1) (0.5714285714,1) (0.5892857143,1) (0.6071428571,1) (0.625,1) (0.6428571429,1) (0.6607142857,1) (0.6785714286,1) (0.6964285714,1) (0.7142857143,1) (0.7321428571,1) (0.75,1) (0.7678571429,1) (0.7857142857,1) (0.8035714286,1) (0.8214285714,1) (0.8392857143,1) (0.8571428571,1) (0.875,1) (0.8928571429,1) (0.9107142857,1) (0.9285714286,1) (0.9464285714,1) (0.9642857143,1) (0.9821428571,1) (1,1)};
\addlegendentry{Outer-fold curves}
\addplot[forget plot,evalgray,line width=0.6pt] coordinates {(0,0) (0,0.1666666667) (0,0.3333333333) (0,0.5) (0.01818181818,0.5) (0.03636363636,0.5) (0.05454545455,0.5) (0.07272727273,0.5) (0.07272727273,0.6666666667) (0.07272727273,0.8333333333) (0.09090909091,0.8333333333) (0.1090909091,0.8333333333) (0.1272727273,0.8333333333) (0.1454545455,0.8333333333) (0.1636363636,0.8333333333) (0.1818181818,0.8333333333) (0.2,0.8333333333) (0.2181818182,0.8333333333) (0.2363636364,0.8333333333) (0.2545454545,0.8333333333) (0.2727272727,0.8333333333) (0.2909090909,0.8333333333) (0.3090909091,0.8333333333) (0.3272727273,0.8333333333) (0.3454545455,0.8333333333) (0.3636363636,0.8333333333) (0.3818181818,0.8333333333) (0.4,0.8333333333) (0.4181818182,0.8333333333) (0.4363636364,0.8333333333) (0.4545454545,0.8333333333) (0.4727272727,0.8333333333) (0.4909090909,0.8333333333) (0.5090909091,0.8333333333) (0.5272727273,0.8333333333) (0.5454545455,0.8333333333) (0.5636363636,0.8333333333) (0.5818181818,0.8333333333) (0.6,0.8333333333) (0.6181818182,0.8333333333) (0.6363636364,0.8333333333) (0.6545454545,0.8333333333) (0.6727272727,0.8333333333) (0.6909090909,0.8333333333) (0.7090909091,0.8333333333) (0.7272727273,0.8333333333) (0.7454545455,0.8333333333) (0.7636363636,0.8333333333) (0.7818181818,0.8333333333) (0.8,0.8333333333) (0.8,1) (0.8181818182,1) (0.8363636364,1) (0.8545454545,1) (0.8727272727,1) (0.8909090909,1) (0.9090909091,1) (0.9272727273,1) (0.9454545455,1) (0.9636363636,1) (0.9818181818,1) (1,1)};
\addplot[forget plot,evalgray,line width=0.6pt] coordinates {(0,0) (0,0.1666666667) (0.01818181818,0.1666666667) (0.01818181818,0.3333333333) (0.01818181818,0.5) (0.03636363636,0.5) (0.03636363636,0.6666666667) (0.03636363636,0.8333333333) (0.05454545455,0.8333333333) (0.07272727273,0.8333333333) (0.09090909091,0.8333333333) (0.1090909091,0.8333333333) (0.1272727273,0.8333333333) (0.1454545455,0.8333333333) (0.1636363636,0.8333333333) (0.1818181818,0.8333333333) (0.2,0.8333333333) (0.2181818182,0.8333333333) (0.2363636364,0.8333333333) (0.2545454545,0.8333333333) (0.2727272727,0.8333333333) (0.2909090909,0.8333333333) (0.3090909091,0.8333333333) (0.3272727273,0.8333333333) (0.3272727273,1) (0.3454545455,1) (0.3636363636,1) (0.3818181818,1) (0.4,1) (0.4181818182,1) (0.4363636364,1) (0.4545454545,1) (0.4727272727,1) (0.4909090909,1) (0.5090909091,1) (0.5272727273,1) (0.5454545455,1) (0.5636363636,1) (0.5818181818,1) (0.6,1) (0.6181818182,1) (0.6363636364,1) (0.6545454545,1) (0.6727272727,1) (0.6909090909,1) (0.7090909091,1) (0.7272727273,1) (0.7454545455,1) (0.7636363636,1) (0.7818181818,1) (0.8,1) (0.8181818182,1) (0.8363636364,1) (0.8545454545,1) (0.8727272727,1) (0.8909090909,1) (0.9090909091,1) (0.9272727273,1) (0.9454545455,1) (0.9636363636,1) (0.9818181818,1) (1,1)};
\addplot[forget plot,evalgray,line width=0.6pt] coordinates {(0,0) (0,0.1666666667) (0,0.3333333333) (0,0.5) (0,0.6666666667) (0,0.8333333333) (0,1) (0.01818181818,1) (0.03636363636,1) (0.05454545455,1) (0.07272727273,1) (0.09090909091,1) (0.1090909091,1) (0.1272727273,1) (0.1454545455,1) (0.1636363636,1) (0.1818181818,1) (0.2,1) (0.2181818182,1) (0.2363636364,1) (0.2545454545,1) (0.2727272727,1) (0.2909090909,1) (0.3090909091,1) (0.3272727273,1) (0.3454545455,1) (0.3636363636,1) (0.3818181818,1) (0.4,1) (0.4181818182,1) (0.4363636364,1) (0.4545454545,1) (0.4727272727,1) (0.4909090909,1) (0.5090909091,1) (0.5272727273,1) (0.5454545455,1) (0.5636363636,1) (0.5818181818,1) (0.6,1) (0.6181818182,1) (0.6363636364,1) (0.6545454545,1) (0.6727272727,1) (0.6909090909,1) (0.7090909091,1) (0.7272727273,1) (0.7454545455,1) (0.7636363636,1) (0.7818181818,1) (0.8,1) (0.8181818182,1) (0.8363636364,1) (0.8545454545,1) (0.8727272727,1) (0.8909090909,1) (0.9090909091,1) (0.9272727273,1) (0.9454545455,1) (0.9636363636,1) (0.9818181818,1) (1,1)};
\addplot[forget plot,evalgray,line width=0.6pt] coordinates {(0,0) (0.01818181818,0) (0.03636363636,0) (0.03636363636,0.2) (0.05454545455,0.2) (0.05454545455,0.4) (0.07272727273,0.4) (0.09090909091,0.4) (0.1090909091,0.4) (0.1272727273,0.4) (0.1454545455,0.4) (0.1636363636,0.4) (0.1818181818,0.4) (0.2,0.4) (0.2,0.6) (0.2181818182,0.6) (0.2363636364,0.6) (0.2545454545,0.6) (0.2727272727,0.6) (0.2909090909,0.6) (0.3090909091,0.6) (0.3272727273,0.6) (0.3454545455,0.6) (0.3454545455,0.8) (0.3636363636,0.8) (0.3818181818,0.8) (0.3818181818,1) (0.4,1) (0.4181818182,1) (0.4363636364,1) (0.4545454545,1) (0.4727272727,1) (0.4909090909,1) (0.5090909091,1) (0.5272727273,1) (0.5454545455,1) (0.5636363636,1) (0.5818181818,1) (0.6,1) (0.6181818182,1) (0.6363636364,1) (0.6545454545,1) (0.6727272727,1) (0.6909090909,1) (0.7090909091,1) (0.7272727273,1) (0.7454545455,1) (0.7636363636,1) (0.7818181818,1) (0.8,1) (0.8181818182,1) (0.8363636364,1) (0.8545454545,1) (0.8727272727,1) (0.8909090909,1) (0.9090909091,1) (0.9272727273,1) (0.9454545455,1) (0.9636363636,1) (0.9818181818,1) (1,1)};
\addplot[evalteal,line width=1.35pt] coordinates {(0,0) (0.003623188406,0) (0.003623188406,0.03571428571) (0.003623188406,0.07142857143) (0.003623188406,0.1071428571) (0.007246376812,0.1071428571) (0.007246376812,0.1428571429) (0.01086956522,0.1428571429) (0.01086956522,0.1785714286) (0.01086956522,0.2142857143) (0.01086956522,0.25) (0.01086956522,0.2857142857) (0.01449275362,0.2857142857) (0.01449275362,0.3214285714) (0.01449275362,0.3571428571) (0.01449275362,0.3928571429) (0.01449275362,0.4285714286) (0.01449275362,0.4642857143) (0.01449275362,0.5) (0.01811594203,0.5) (0.02173913043,0.5) (0.02536231884,0.5) (0.02536231884,0.5357142857) (0.02536231884,0.5714285714) (0.02536231884,0.6071428571) (0.02536231884,0.6428571429) (0.02536231884,0.6785714286) (0.02898550725,0.6785714286) (0.03260869565,0.6785714286) (0.03623188406,0.6785714286) (0.03985507246,0.6785714286) (0.04347826087,0.6785714286) (0.04710144928,0.6785714286) (0.05072463768,0.6785714286) (0.05434782609,0.6785714286) (0.05797101449,0.6785714286) (0.0615942029,0.6785714286) (0.0652173913,0.6785714286) (0.0652173913,0.7142857143) (0.0652173913,0.75) (0.06884057971,0.75) (0.07246376812,0.75) (0.07608695652,0.75) (0.07971014493,0.75) (0.08333333333,0.75) (0.08695652174,0.75) (0.08695652174,0.7857142857) (0.09057971014,0.7857142857) (0.09420289855,0.7857142857) (0.09782608696,0.7857142857) (0.1014492754,0.7857142857) (0.1050724638,0.7857142857) (0.1086956522,0.7857142857) (0.1123188406,0.7857142857) (0.115942029,0.7857142857) (0.1195652174,0.7857142857) (0.1231884058,0.7857142857) (0.1268115942,0.7857142857) (0.1304347826,0.7857142857) (0.134057971,0.7857142857) (0.134057971,0.8214285714) (0.1376811594,0.8214285714) (0.1413043478,0.8214285714) (0.1449275362,0.8214285714) (0.1485507246,0.8214285714) (0.152173913,0.8214285714) (0.1557971014,0.8214285714) (0.1594202899,0.8214285714) (0.1630434783,0.8214285714) (0.1666666667,0.8214285714) (0.1702898551,0.8214285714) (0.1739130435,0.8214285714) (0.1775362319,0.8214285714) (0.1811594203,0.8214285714) (0.1847826087,0.8214285714) (0.1884057971,0.8214285714) (0.1920289855,0.8214285714) (0.1956521739,0.8214285714) (0.1992753623,0.8214285714) (0.2028985507,0.8214285714) (0.2065217391,0.8214285714) (0.2101449275,0.8214285714) (0.2137681159,0.8214285714) (0.2173913043,0.8214285714) (0.2210144928,0.8214285714) (0.2246376812,0.8214285714) (0.2282608696,0.8214285714) (0.231884058,0.8214285714) (0.2355072464,0.8214285714) (0.2391304348,0.8214285714) (0.2427536232,0.8214285714) (0.2463768116,0.8214285714) (0.25,0.8214285714) (0.25,0.8571428571) (0.2536231884,0.8571428571) (0.2572463768,0.8571428571) (0.2608695652,0.8571428571) (0.2644927536,0.8571428571) (0.268115942,0.8571428571) (0.2717391304,0.8571428571) (0.2753623188,0.8571428571) (0.2789855072,0.8571428571) (0.2826086957,0.8571428571) (0.2862318841,0.8571428571) (0.2898550725,0.8571428571) (0.2934782609,0.8571428571) (0.2971014493,0.8571428571) (0.3007246377,0.8571428571) (0.3043478261,0.8571428571) (0.3079710145,0.8571428571) (0.3115942029,0.8571428571) (0.3115942029,0.8928571429) (0.3152173913,0.8928571429) (0.3188405797,0.8928571429) (0.3224637681,0.8928571429) (0.3260869565,0.8928571429) (0.3297101449,0.8928571429) (0.3333333333,0.8928571429) (0.3369565217,0.8928571429) (0.3405797101,0.8928571429) (0.3442028986,0.8928571429) (0.347826087,0.8928571429) (0.3514492754,0.8928571429) (0.3550724638,0.8928571429) (0.3586956522,0.8928571429) (0.3623188406,0.8928571429) (0.365942029,0.8928571429) (0.3695652174,0.8928571429) (0.3731884058,0.8928571429) (0.3768115942,0.8928571429) (0.3804347826,0.8928571429) (0.384057971,0.8928571429) (0.3876811594,0.8928571429) (0.3913043478,0.8928571429) (0.3949275362,0.8928571429) (0.3985507246,0.8928571429) (0.402173913,0.8928571429) (0.402173913,0.9285714286) (0.4057971014,0.9285714286) (0.4094202899,0.9285714286) (0.4130434783,0.9285714286) (0.4166666667,0.9285714286) (0.4202898551,0.9285714286) (0.4202898551,0.9642857143) (0.4239130435,0.9642857143) (0.4275362319,0.9642857143) (0.4311594203,0.9642857143) (0.4347826087,0.9642857143) (0.4384057971,0.9642857143) (0.4420289855,0.9642857143) (0.4456521739,0.9642857143) (0.4492753623,0.9642857143) (0.4528985507,0.9642857143) (0.4565217391,0.9642857143) (0.4601449275,0.9642857143) (0.4637681159,0.9642857143) (0.4673913043,0.9642857143) (0.4710144928,0.9642857143) (0.4746376812,0.9642857143) (0.4782608696,0.9642857143) (0.481884058,0.9642857143) (0.4855072464,0.9642857143) (0.4891304348,0.9642857143) (0.4927536232,0.9642857143) (0.4963768116,0.9642857143) (0.5,0.9642857143) (0.5036231884,0.9642857143) (0.5072463768,0.9642857143) (0.5108695652,0.9642857143) (0.5144927536,0.9642857143) (0.518115942,0.9642857143) (0.5217391304,0.9642857143) (0.5253623188,0.9642857143) (0.5289855072,0.9642857143) (0.5326086957,0.9642857143) (0.5362318841,0.9642857143) (0.5398550725,0.9642857143) (0.5434782609,0.9642857143) (0.5471014493,0.9642857143) (0.5507246377,0.9642857143) (0.5543478261,0.9642857143) (0.5579710145,0.9642857143) (0.5615942029,0.9642857143) (0.5652173913,0.9642857143) (0.5688405797,0.9642857143) (0.5724637681,0.9642857143) (0.5760869565,0.9642857143) (0.5797101449,0.9642857143) (0.5833333333,0.9642857143) (0.5869565217,0.9642857143) (0.5905797101,0.9642857143) (0.5942028986,0.9642857143) (0.597826087,0.9642857143) (0.6014492754,0.9642857143) (0.6050724638,0.9642857143) (0.6086956522,0.9642857143) (0.6123188406,0.9642857143) (0.615942029,0.9642857143) (0.6195652174,0.9642857143) (0.6231884058,0.9642857143) (0.6268115942,0.9642857143) (0.6304347826,0.9642857143) (0.634057971,0.9642857143) (0.6376811594,0.9642857143) (0.6413043478,0.9642857143) (0.6449275362,0.9642857143) (0.6485507246,0.9642857143) (0.652173913,0.9642857143) (0.6557971014,0.9642857143) (0.6594202899,0.9642857143) (0.6630434783,0.9642857143) (0.6666666667,0.9642857143) (0.6702898551,0.9642857143) (0.6739130435,0.9642857143) (0.6775362319,0.9642857143) (0.6811594203,0.9642857143) (0.6847826087,0.9642857143) (0.6884057971,0.9642857143) (0.6920289855,0.9642857143) (0.6956521739,0.9642857143) (0.6992753623,0.9642857143) (0.7028985507,0.9642857143) (0.7065217391,0.9642857143) (0.7101449275,0.9642857143) (0.7137681159,0.9642857143) (0.7173913043,0.9642857143) (0.7210144928,0.9642857143) (0.7246376812,0.9642857143) (0.7282608696,0.9642857143) (0.731884058,0.9642857143) (0.7355072464,0.9642857143) (0.7391304348,0.9642857143) (0.7427536232,0.9642857143) (0.7463768116,0.9642857143) (0.75,0.9642857143) (0.7536231884,0.9642857143) (0.7572463768,0.9642857143) (0.7608695652,0.9642857143) (0.7644927536,0.9642857143) (0.768115942,0.9642857143) (0.7717391304,0.9642857143) (0.7753623188,0.9642857143) (0.7789855072,0.9642857143) (0.7826086957,0.9642857143) (0.7862318841,0.9642857143) (0.7898550725,0.9642857143) (0.7898550725,1) (0.7934782609,1) (0.7971014493,1) (0.8007246377,1) (0.8043478261,1) (0.8079710145,1) (0.8115942029,1) (0.8152173913,1) (0.8188405797,1) (0.8224637681,1) (0.8260869565,1) (0.8297101449,1) (0.8333333333,1) (0.8369565217,1) (0.8405797101,1) (0.8442028986,1) (0.847826087,1) (0.8514492754,1) (0.8550724638,1) (0.8586956522,1) (0.8623188406,1) (0.865942029,1) (0.8695652174,1) (0.8731884058,1) (0.8768115942,1) (0.8804347826,1) (0.884057971,1) (0.8876811594,1) (0.8913043478,1) (0.8949275362,1) (0.8985507246,1) (0.902173913,1) (0.9057971014,1) (0.9094202899,1) (0.9130434783,1) (0.9166666667,1) (0.9202898551,1) (0.9239130435,1) (0.9275362319,1) (0.9311594203,1) (0.9347826087,1) (0.9384057971,1) (0.9420289855,1) (0.9456521739,1) (0.9492753623,1) (0.9528985507,1) (0.9565217391,1) (0.9601449275,1) (0.9637681159,1) (0.9673913043,1) (0.9710144928,1) (0.9746376812,1) (0.9782608696,1) (0.981884058,1) (0.9855072464,1) (0.9891304348,1) (0.9927536232,1) (0.9963768116,1) (1,1)};
\addlegendentry{Pooled AUROC = 0.900}
\addplot[evalorange!70,dashed,line width=0.7pt] coordinates {(0,0) (1,1)};
\addlegendentry{Chance diagonal}
\addplot[forget plot,only marks,mark=*,mark size=1.8pt,evalorange] coordinates {(0.01449275362,0.3571428571)};
\node[anchor=west,font=\scriptsize,text=evalorange] at (axis cs:0.069492754,0.28214286) {$s>0$};
\draw[evalorange,thin,->] (axis cs:0.079492754,0.32714286) -- (axis cs:0.014492754,0.35714286);
\nextgroupplot[title={(b) ET precision--recall},
 xlabel={Recall (ET)},ylabel={Precision (ET)},
 legend style={at={(0.5,-0.30)},anchor=north}]
\addplot[evalgray,line width=0.6pt] coordinates {(0,1) (0,1) (0.2,1) (0.2,0.5) (0.2,0.5) (0.2,0.6666666667) (0.4,0.6666666667) (0.4,0.75) (0.6,0.75) (0.6,0.6) (0.6,0.6) (0.6,0.5) (0.6,0.5) (0.6,0.5714285714) (0.8,0.5714285714) (0.8,0.5) (0.8,0.5) (0.8,0.5555555556) (1,0.5555555556) (1,0.5) (1,0.5) (1,0.4545454545) (1,0.4545454545) (1,0.4166666667) (1,0.4166666667) (1,0.3846153846) (1,0.3846153846) (1,0.3571428571) (1,0.3571428571) (1,0.3333333333) (1,0.3333333333) (1,0.3125) (1,0.3125) (1,0.2941176471) (1,0.2941176471) (1,0.2777777778) (1,0.2777777778) (1,0.2631578947) (1,0.2631578947) (1,0.25) (1,0.25) (1,0.2380952381) (1,0.2380952381) (1,0.2272727273) (1,0.2272727273) (1,0.2173913043) (1,0.2173913043) (1,0.2083333333) (1,0.2083333333) (1,0.2) (1,0.2) (1,0.1923076923) (1,0.1923076923) (1,0.1851851852) (1,0.1851851852) (1,0.1785714286) (1,0.1785714286) (1,0.1724137931) (1,0.1724137931) (1,0.1666666667) (1,0.1666666667) (1,0.1612903226) (1,0.1612903226) (1,0.15625) (1,0.15625) (1,0.1515151515) (1,0.1515151515) (1,0.1470588235) (1,0.1470588235) (1,0.1428571429) (1,0.1428571429) (1,0.1388888889) (1,0.1388888889) (1,0.1351351351) (1,0.1351351351) (1,0.1315789474) (1,0.1315789474) (1,0.1282051282) (1,0.1282051282) (1,0.125) (1,0.125) (1,0.1219512195) (1,0.1219512195) (1,0.119047619) (1,0.119047619) (1,0.1162790698) (1,0.1162790698) (1,0.1136363636) (1,0.1136363636) (1,0.1111111111) (1,0.1111111111) (1,0.1086956522) (1,0.1086956522) (1,0.1063829787) (1,0.1063829787) (1,0.1041666667) (1,0.1041666667) (1,0.1020408163) (1,0.1020408163) (1,0.1) (1,0.1) (1,0.09803921569) (1,0.09803921569) (1,0.09615384615) (1,0.09615384615) (1,0.09433962264) (1,0.09433962264) (1,0.09259259259) (1,0.09259259259) (1,0.09090909091) (1,0.09090909091) (1,0.08928571429) (1,0.08928571429) (1,0.08771929825) (1,0.08771929825) (1,0.08620689655) (1,0.08620689655) (1,0.08474576271) (1,0.08474576271) (1,0.08333333333) (1,0.08333333333) (1,0.08196721311) (1,0.08196721311)};
\addlegendentry{Outer-fold curves}
\addplot[forget plot,evalgray,line width=0.6pt] coordinates {(0,1) (0,1) (0.1666666667,1) (0.1666666667,1) (0.3333333333,1) (0.3333333333,1) (0.5,1) (0.5,0.75) (0.5,0.75) (0.5,0.6) (0.5,0.6) (0.5,0.5) (0.5,0.5) (0.5,0.4285714286) (0.5,0.4285714286) (0.5,0.5) (0.6666666667,0.5) (0.6666666667,0.5555555556) (0.8333333333,0.5555555556) (0.8333333333,0.5) (0.8333333333,0.5) (0.8333333333,0.4545454545) (0.8333333333,0.4545454545) (0.8333333333,0.4166666667) (0.8333333333,0.4166666667) (0.8333333333,0.3846153846) (0.8333333333,0.3846153846) (0.8333333333,0.3571428571) (0.8333333333,0.3571428571) (0.8333333333,0.3333333333) (0.8333333333,0.3333333333) (0.8333333333,0.3125) (0.8333333333,0.3125) (0.8333333333,0.2941176471) (0.8333333333,0.2941176471) (0.8333333333,0.2777777778) (0.8333333333,0.2777777778) (0.8333333333,0.2631578947) (0.8333333333,0.2631578947) (0.8333333333,0.25) (0.8333333333,0.25) (0.8333333333,0.2380952381) (0.8333333333,0.2380952381) (0.8333333333,0.2272727273) (0.8333333333,0.2272727273) (0.8333333333,0.2173913043) (0.8333333333,0.2173913043) (0.8333333333,0.2083333333) (0.8333333333,0.2083333333) (0.8333333333,0.2) (0.8333333333,0.2) (0.8333333333,0.1923076923) (0.8333333333,0.1923076923) (0.8333333333,0.1851851852) (0.8333333333,0.1851851852) (0.8333333333,0.1785714286) (0.8333333333,0.1785714286) (0.8333333333,0.1724137931) (0.8333333333,0.1724137931) (0.8333333333,0.1666666667) (0.8333333333,0.1666666667) (0.8333333333,0.1612903226) (0.8333333333,0.1612903226) (0.8333333333,0.15625) (0.8333333333,0.15625) (0.8333333333,0.1515151515) (0.8333333333,0.1515151515) (0.8333333333,0.1470588235) (0.8333333333,0.1470588235) (0.8333333333,0.1428571429) (0.8333333333,0.1428571429) (0.8333333333,0.1388888889) (0.8333333333,0.1388888889) (0.8333333333,0.1351351351) (0.8333333333,0.1351351351) (0.8333333333,0.1315789474) (0.8333333333,0.1315789474) (0.8333333333,0.1282051282) (0.8333333333,0.1282051282) (0.8333333333,0.125) (0.8333333333,0.125) (0.8333333333,0.1219512195) (0.8333333333,0.1219512195) (0.8333333333,0.119047619) (0.8333333333,0.119047619) (0.8333333333,0.1162790698) (0.8333333333,0.1162790698) (0.8333333333,0.1136363636) (0.8333333333,0.1136363636) (0.8333333333,0.1111111111) (0.8333333333,0.1111111111) (0.8333333333,0.1086956522) (0.8333333333,0.1086956522) (0.8333333333,0.1063829787) (0.8333333333,0.1063829787) (0.8333333333,0.1041666667) (0.8333333333,0.1041666667) (0.8333333333,0.1020408163) (0.8333333333,0.1020408163) (0.8333333333,0.12) (1,0.12) (1,0.1176470588) (1,0.1176470588) (1,0.1153846154) (1,0.1153846154) (1,0.1132075472) (1,0.1132075472) (1,0.1111111111) (1,0.1111111111) (1,0.1090909091) (1,0.1090909091) (1,0.1071428571) (1,0.1071428571) (1,0.1052631579) (1,0.1052631579) (1,0.1034482759) (1,0.1034482759) (1,0.1016949153) (1,0.1016949153) (1,0.1) (1,0.1) (1,0.09836065574) (1,0.09836065574)};
\addplot[forget plot,evalgray,line width=0.6pt] coordinates {(0,1) (0,1) (0.1666666667,1) (0.1666666667,0.5) (0.1666666667,0.5) (0.1666666667,0.6666666667) (0.3333333333,0.6666666667) (0.3333333333,0.75) (0.5,0.75) (0.5,0.6) (0.5,0.6) (0.5,0.6666666667) (0.6666666667,0.6666666667) (0.6666666667,0.7142857143) (0.8333333333,0.7142857143) (0.8333333333,0.625) (0.8333333333,0.625) (0.8333333333,0.5555555556) (0.8333333333,0.5555555556) (0.8333333333,0.5) (0.8333333333,0.5) (0.8333333333,0.4545454545) (0.8333333333,0.4545454545) (0.8333333333,0.4166666667) (0.8333333333,0.4166666667) (0.8333333333,0.3846153846) (0.8333333333,0.3846153846) (0.8333333333,0.3571428571) (0.8333333333,0.3571428571) (0.8333333333,0.3333333333) (0.8333333333,0.3333333333) (0.8333333333,0.3125) (0.8333333333,0.3125) (0.8333333333,0.2941176471) (0.8333333333,0.2941176471) (0.8333333333,0.2777777778) (0.8333333333,0.2777777778) (0.8333333333,0.2631578947) (0.8333333333,0.2631578947) (0.8333333333,0.25) (0.8333333333,0.25) (0.8333333333,0.2380952381) (0.8333333333,0.2380952381) (0.8333333333,0.2272727273) (0.8333333333,0.2272727273) (0.8333333333,0.2173913043) (0.8333333333,0.2173913043) (0.8333333333,0.25) (1,0.25) (1,0.24) (1,0.24) (1,0.2307692308) (1,0.2307692308) (1,0.2222222222) (1,0.2222222222) (1,0.2142857143) (1,0.2142857143) (1,0.2068965517) (1,0.2068965517) (1,0.2) (1,0.2) (1,0.1935483871) (1,0.1935483871) (1,0.1875) (1,0.1875) (1,0.1818181818) (1,0.1818181818) (1,0.1764705882) (1,0.1764705882) (1,0.1714285714) (1,0.1714285714) (1,0.1666666667) (1,0.1666666667) (1,0.1621621622) (1,0.1621621622) (1,0.1578947368) (1,0.1578947368) (1,0.1538461538) (1,0.1538461538) (1,0.15) (1,0.15) (1,0.1463414634) (1,0.1463414634) (1,0.1428571429) (1,0.1428571429) (1,0.1395348837) (1,0.1395348837) (1,0.1363636364) (1,0.1363636364) (1,0.1333333333) (1,0.1333333333) (1,0.1304347826) (1,0.1304347826) (1,0.1276595745) (1,0.1276595745) (1,0.125) (1,0.125) (1,0.1224489796) (1,0.1224489796) (1,0.12) (1,0.12) (1,0.1176470588) (1,0.1176470588) (1,0.1153846154) (1,0.1153846154) (1,0.1132075472) (1,0.1132075472) (1,0.1111111111) (1,0.1111111111) (1,0.1090909091) (1,0.1090909091) (1,0.1071428571) (1,0.1071428571) (1,0.1052631579) (1,0.1052631579) (1,0.1034482759) (1,0.1034482759) (1,0.1016949153) (1,0.1016949153) (1,0.1) (1,0.1) (1,0.09836065574) (1,0.09836065574)};
\addplot[forget plot,evalgray,line width=0.6pt] coordinates {(0,1) (0,1) (0.1666666667,1) (0.1666666667,1) (0.3333333333,1) (0.3333333333,1) (0.5,1) (0.5,1) (0.6666666667,1) (0.6666666667,1) (0.8333333333,1) (0.8333333333,1) (1,1) (1,0.8571428571) (1,0.8571428571) (1,0.75) (1,0.75) (1,0.6666666667) (1,0.6666666667) (1,0.6) (1,0.6) (1,0.5454545455) (1,0.5454545455) (1,0.5) (1,0.5) (1,0.4615384615) (1,0.4615384615) (1,0.4285714286) (1,0.4285714286) (1,0.4) (1,0.4) (1,0.375) (1,0.375) (1,0.3529411765) (1,0.3529411765) (1,0.3333333333) (1,0.3333333333) (1,0.3157894737) (1,0.3157894737) (1,0.3) (1,0.3) (1,0.2857142857) (1,0.2857142857) (1,0.2727272727) (1,0.2727272727) (1,0.2608695652) (1,0.2608695652) (1,0.25) (1,0.25) (1,0.24) (1,0.24) (1,0.2307692308) (1,0.2307692308) (1,0.2222222222) (1,0.2222222222) (1,0.2142857143) (1,0.2142857143) (1,0.2068965517) (1,0.2068965517) (1,0.2) (1,0.2) (1,0.1935483871) (1,0.1935483871) (1,0.1875) (1,0.1875) (1,0.1818181818) (1,0.1818181818) (1,0.1764705882) (1,0.1764705882) (1,0.1714285714) (1,0.1714285714) (1,0.1666666667) (1,0.1666666667) (1,0.1621621622) (1,0.1621621622) (1,0.1578947368) (1,0.1578947368) (1,0.1538461538) (1,0.1538461538) (1,0.15) (1,0.15) (1,0.1463414634) (1,0.1463414634) (1,0.1428571429) (1,0.1428571429) (1,0.1395348837) (1,0.1395348837) (1,0.1363636364) (1,0.1363636364) (1,0.1333333333) (1,0.1333333333) (1,0.1304347826) (1,0.1304347826) (1,0.1276595745) (1,0.1276595745) (1,0.125) (1,0.125) (1,0.1224489796) (1,0.1224489796) (1,0.12) (1,0.12) (1,0.1176470588) (1,0.1176470588) (1,0.1153846154) (1,0.1153846154) (1,0.1132075472) (1,0.1132075472) (1,0.1111111111) (1,0.1111111111) (1,0.1090909091) (1,0.1090909091) (1,0.1071428571) (1,0.1071428571) (1,0.1052631579) (1,0.1052631579) (1,0.1034482759) (1,0.1034482759) (1,0.1016949153) (1,0.1016949153) (1,0.1) (1,0.1) (1,0.09836065574) (1,0.09836065574)};
\addplot[forget plot,evalgray,line width=0.6pt] coordinates {(0,1) (0,0) (0,0) (0,0) (0,0) (0,0.3333333333) (0.2,0.3333333333) (0.2,0.25) (0.2,0.25) (0.2,0.4) (0.4,0.4) (0.4,0.3333333333) (0.4,0.3333333333) (0.4,0.2857142857) (0.4,0.2857142857) (0.4,0.25) (0.4,0.25) (0.4,0.2222222222) (0.4,0.2222222222) (0.4,0.2) (0.4,0.2) (0.4,0.1818181818) (0.4,0.1818181818) (0.4,0.1666666667) (0.4,0.1666666667) (0.4,0.1538461538) (0.4,0.1538461538) (0.4,0.2142857143) (0.6,0.2142857143) (0.6,0.2) (0.6,0.2) (0.6,0.1875) (0.6,0.1875) (0.6,0.1764705882) (0.6,0.1764705882) (0.6,0.1666666667) (0.6,0.1666666667) (0.6,0.1578947368) (0.6,0.1578947368) (0.6,0.15) (0.6,0.15) (0.6,0.1428571429) (0.6,0.1428571429) (0.6,0.1363636364) (0.6,0.1363636364) (0.6,0.1739130435) (0.8,0.1739130435) (0.8,0.1666666667) (0.8,0.1666666667) (0.8,0.16) (0.8,0.16) (0.8,0.1923076923) (1,0.1923076923) (1,0.1851851852) (1,0.1851851852) (1,0.1785714286) (1,0.1785714286) (1,0.1724137931) (1,0.1724137931) (1,0.1666666667) (1,0.1666666667) (1,0.1612903226) (1,0.1612903226) (1,0.15625) (1,0.15625) (1,0.1515151515) (1,0.1515151515) (1,0.1470588235) (1,0.1470588235) (1,0.1428571429) (1,0.1428571429) (1,0.1388888889) (1,0.1388888889) (1,0.1351351351) (1,0.1351351351) (1,0.1315789474) (1,0.1315789474) (1,0.1282051282) (1,0.1282051282) (1,0.125) (1,0.125) (1,0.1219512195) (1,0.1219512195) (1,0.119047619) (1,0.119047619) (1,0.1162790698) (1,0.1162790698) (1,0.1136363636) (1,0.1136363636) (1,0.1111111111) (1,0.1111111111) (1,0.1086956522) (1,0.1086956522) (1,0.1063829787) (1,0.1063829787) (1,0.1041666667) (1,0.1041666667) (1,0.1020408163) (1,0.1020408163) (1,0.1) (1,0.1) (1,0.09803921569) (1,0.09803921569) (1,0.09615384615) (1,0.09615384615) (1,0.09433962264) (1,0.09433962264) (1,0.09259259259) (1,0.09259259259) (1,0.09090909091) (1,0.09090909091) (1,0.08928571429) (1,0.08928571429) (1,0.08771929825) (1,0.08771929825) (1,0.08620689655) (1,0.08620689655) (1,0.08474576271) (1,0.08474576271) (1,0.08333333333) (1,0.08333333333)};
\addplot[evalteal,line width=1.35pt] coordinates {(0,1) (0,0) (0,0) (0,0.5) (0.03571428571,0.5) (0.03571428571,0.6666666667) (0.07142857143,0.6666666667) (0.07142857143,0.75) (0.1071428571,0.75) (0.1071428571,0.6) (0.1071428571,0.6) (0.1071428571,0.6666666667) (0.1428571429,0.6666666667) (0.1428571429,0.5714285714) (0.1428571429,0.5714285714) (0.1428571429,0.625) (0.1785714286,0.625) (0.1785714286,0.6666666667) (0.2142857143,0.6666666667) (0.2142857143,0.7) (0.25,0.7) (0.25,0.7272727273) (0.2857142857,0.7272727273) (0.2857142857,0.6666666667) (0.2857142857,0.6666666667) (0.2857142857,0.6923076923) (0.3214285714,0.6923076923) (0.3214285714,0.7142857143) (0.3571428571,0.7142857143) (0.3571428571,0.7333333333) (0.3928571429,0.7333333333) (0.3928571429,0.75) (0.4285714286,0.75) (0.4285714286,0.7647058824) (0.4642857143,0.7647058824) (0.4642857143,0.7777777778) (0.5,0.7777777778) (0.5,0.7368421053) (0.5,0.7368421053) (0.5,0.7) (0.5,0.7) (0.5,0.6666666667) (0.5,0.6666666667) (0.5,0.6818181818) (0.5357142857,0.6818181818) (0.5357142857,0.6956521739) (0.5714285714,0.6956521739) (0.5714285714,0.7083333333) (0.6071428571,0.7083333333) (0.6071428571,0.72) (0.6428571429,0.72) (0.6428571429,0.7307692308) (0.6785714286,0.7307692308) (0.6785714286,0.7037037037) (0.6785714286,0.7037037037) (0.6785714286,0.6785714286) (0.6785714286,0.6785714286) (0.6785714286,0.6551724138) (0.6785714286,0.6551724138) (0.6785714286,0.6333333333) (0.6785714286,0.6333333333) (0.6785714286,0.6129032258) (0.6785714286,0.6129032258) (0.6785714286,0.59375) (0.6785714286,0.59375) (0.6785714286,0.5757575758) (0.6785714286,0.5757575758) (0.6785714286,0.5588235294) (0.6785714286,0.5588235294) (0.6785714286,0.5428571429) (0.6785714286,0.5428571429) (0.6785714286,0.5277777778) (0.6785714286,0.5277777778) (0.6785714286,0.5135135135) (0.6785714286,0.5135135135) (0.6785714286,0.5263157895) (0.7142857143,0.5263157895) (0.7142857143,0.5384615385) (0.75,0.5384615385) (0.75,0.525) (0.75,0.525) (0.75,0.512195122) (0.75,0.512195122) (0.75,0.5) (0.75,0.5) (0.75,0.488372093) (0.75,0.488372093) (0.75,0.4772727273) (0.75,0.4772727273) (0.75,0.4666666667) (0.75,0.4666666667) (0.75,0.4782608696) (0.7857142857,0.4782608696) (0.7857142857,0.4680851064) (0.7857142857,0.4680851064) (0.7857142857,0.4583333333) (0.7857142857,0.4583333333) (0.7857142857,0.4489795918) (0.7857142857,0.4489795918) (0.7857142857,0.44) (0.7857142857,0.44) (0.7857142857,0.431372549) (0.7857142857,0.431372549) (0.7857142857,0.4230769231) (0.7857142857,0.4230769231) (0.7857142857,0.4150943396) (0.7857142857,0.4150943396) (0.7857142857,0.4074074074) (0.7857142857,0.4074074074) (0.7857142857,0.4) (0.7857142857,0.4) (0.7857142857,0.3928571429) (0.7857142857,0.3928571429) (0.7857142857,0.3859649123) (0.7857142857,0.3859649123) (0.7857142857,0.3793103448) (0.7857142857,0.3793103448) (0.7857142857,0.3728813559) (0.7857142857,0.3728813559) (0.7857142857,0.3833333333) (0.8214285714,0.3833333333) (0.8214285714,0.3770491803) (0.8214285714,0.3770491803) (0.8214285714,0.3709677419) (0.8214285714,0.3709677419) (0.8214285714,0.3650793651) (0.8214285714,0.3650793651) (0.8214285714,0.359375) (0.8214285714,0.359375) (0.8214285714,0.3538461538) (0.8214285714,0.3538461538) (0.8214285714,0.3484848485) (0.8214285714,0.3484848485) (0.8214285714,0.3432835821) (0.8214285714,0.3432835821) (0.8214285714,0.3382352941) (0.8214285714,0.3382352941) (0.8214285714,0.3333333333) (0.8214285714,0.3333333333) (0.8214285714,0.3285714286) (0.8214285714,0.3285714286) (0.8214285714,0.323943662) (0.8214285714,0.323943662) (0.8214285714,0.3194444444) (0.8214285714,0.3194444444) (0.8214285714,0.3150684932) (0.8214285714,0.3150684932) (0.8214285714,0.3108108108) (0.8214285714,0.3108108108) (0.8214285714,0.3066666667) (0.8214285714,0.3066666667) (0.8214285714,0.3026315789) (0.8214285714,0.3026315789) (0.8214285714,0.2987012987) (0.8214285714,0.2987012987) (0.8214285714,0.2948717949) (0.8214285714,0.2948717949) (0.8214285714,0.2911392405) (0.8214285714,0.2911392405) (0.8214285714,0.2875) (0.8214285714,0.2875) (0.8214285714,0.2839506173) (0.8214285714,0.2839506173) (0.8214285714,0.2804878049) (0.8214285714,0.2804878049) (0.8214285714,0.2771084337) (0.8214285714,0.2771084337) (0.8214285714,0.2738095238) (0.8214285714,0.2738095238) (0.8214285714,0.2705882353) (0.8214285714,0.2705882353) (0.8214285714,0.2674418605) (0.8214285714,0.2674418605) (0.8214285714,0.2643678161) (0.8214285714,0.2643678161) (0.8214285714,0.2613636364) (0.8214285714,0.2613636364) (0.8214285714,0.2584269663) (0.8214285714,0.2584269663) (0.8214285714,0.2555555556) (0.8214285714,0.2555555556) (0.8214285714,0.2527472527) (0.8214285714,0.2527472527) (0.8214285714,0.25) (0.8214285714,0.25) (0.8214285714,0.2580645161) (0.8571428571,0.2580645161) (0.8571428571,0.2553191489) (0.8571428571,0.2553191489) (0.8571428571,0.2526315789) (0.8571428571,0.2526315789) (0.8571428571,0.25) (0.8571428571,0.25) (0.8571428571,0.2474226804) (0.8571428571,0.2474226804) (0.8571428571,0.2448979592) (0.8571428571,0.2448979592) (0.8571428571,0.2424242424) (0.8571428571,0.2424242424) (0.8571428571,0.24) (0.8571428571,0.24) (0.8571428571,0.2376237624) (0.8571428571,0.2376237624) (0.8571428571,0.2352941176) (0.8571428571,0.2352941176) (0.8571428571,0.2330097087) (0.8571428571,0.2330097087) (0.8571428571,0.2307692308) (0.8571428571,0.2307692308) (0.8571428571,0.2285714286) (0.8571428571,0.2285714286) (0.8571428571,0.2264150943) (0.8571428571,0.2264150943) (0.8571428571,0.2242990654) (0.8571428571,0.2242990654) (0.8571428571,0.2222222222) (0.8571428571,0.2222222222) (0.8571428571,0.2201834862) (0.8571428571,0.2201834862) (0.8571428571,0.2181818182) (0.8571428571,0.2181818182) (0.8571428571,0.2252252252) (0.8928571429,0.2252252252) (0.8928571429,0.2232142857) (0.8928571429,0.2232142857) (0.8928571429,0.2212389381) (0.8928571429,0.2212389381) (0.8928571429,0.2192982456) (0.8928571429,0.2192982456) (0.8928571429,0.2173913043) (0.8928571429,0.2173913043) (0.8928571429,0.2155172414) (0.8928571429,0.2155172414) (0.8928571429,0.2136752137) (0.8928571429,0.2136752137) (0.8928571429,0.2118644068) (0.8928571429,0.2118644068) (0.8928571429,0.2100840336) (0.8928571429,0.2100840336) (0.8928571429,0.2083333333) (0.8928571429,0.2083333333) (0.8928571429,0.2066115702) (0.8928571429,0.2066115702) (0.8928571429,0.2049180328) (0.8928571429,0.2049180328) (0.8928571429,0.2032520325) (0.8928571429,0.2032520325) (0.8928571429,0.2016129032) (0.8928571429,0.2016129032) (0.8928571429,0.2) (0.8928571429,0.2) (0.8928571429,0.1984126984) (0.8928571429,0.1984126984) (0.8928571429,0.1968503937) (0.8928571429,0.1968503937) (0.8928571429,0.1953125) (0.8928571429,0.1953125) (0.8928571429,0.1937984496) (0.8928571429,0.1937984496) (0.8928571429,0.1923076923) (0.8928571429,0.1923076923) (0.8928571429,0.1908396947) (0.8928571429,0.1908396947) (0.8928571429,0.1893939394) (0.8928571429,0.1893939394) (0.8928571429,0.1879699248) (0.8928571429,0.1879699248) (0.8928571429,0.1865671642) (0.8928571429,0.1865671642) (0.8928571429,0.1851851852) (0.8928571429,0.1851851852) (0.8928571429,0.1838235294) (0.8928571429,0.1838235294) (0.8928571429,0.1897810219) (0.9285714286,0.1897810219) (0.9285714286,0.1884057971) (0.9285714286,0.1884057971) (0.9285714286,0.1870503597) (0.9285714286,0.1870503597) (0.9285714286,0.1857142857) (0.9285714286,0.1857142857) (0.9285714286,0.1843971631) (0.9285714286,0.1843971631) (0.9285714286,0.1830985915) (0.9285714286,0.1830985915) (0.9285714286,0.1888111888) (0.9642857143,0.1888111888) (0.9642857143,0.1875) (0.9642857143,0.1875) (0.9642857143,0.1862068966) (0.9642857143,0.1862068966) (0.9642857143,0.1849315068) (0.9642857143,0.1849315068) (0.9642857143,0.1836734694) (0.9642857143,0.1836734694) (0.9642857143,0.1824324324) (0.9642857143,0.1824324324) (0.9642857143,0.1812080537) (0.9642857143,0.1812080537) (0.9642857143,0.18) (0.9642857143,0.18) (0.9642857143,0.178807947) (0.9642857143,0.178807947) (0.9642857143,0.1776315789) (0.9642857143,0.1776315789) (0.9642857143,0.1764705882) (0.9642857143,0.1764705882) (0.9642857143,0.1753246753) (0.9642857143,0.1753246753) (0.9642857143,0.1741935484) (0.9642857143,0.1741935484) (0.9642857143,0.1730769231) (0.9642857143,0.1730769231) (0.9642857143,0.1719745223) (0.9642857143,0.1719745223) (0.9642857143,0.1708860759) (0.9642857143,0.1708860759) (0.9642857143,0.1698113208) (0.9642857143,0.1698113208) (0.9642857143,0.16875) (0.9642857143,0.16875) (0.9642857143,0.1677018634) (0.9642857143,0.1677018634) (0.9642857143,0.1666666667) (0.9642857143,0.1666666667) (0.9642857143,0.1656441718) (0.9642857143,0.1656441718) (0.9642857143,0.1646341463) (0.9642857143,0.1646341463) (0.9642857143,0.1636363636) (0.9642857143,0.1636363636) (0.9642857143,0.1626506024) (0.9642857143,0.1626506024) (0.9642857143,0.1616766467) (0.9642857143,0.1616766467) (0.9642857143,0.1607142857) (0.9642857143,0.1607142857) (0.9642857143,0.1597633136) (0.9642857143,0.1597633136) (0.9642857143,0.1588235294) (0.9642857143,0.1588235294) (0.9642857143,0.1578947368) (0.9642857143,0.1578947368) (0.9642857143,0.1569767442) (0.9642857143,0.1569767442) (0.9642857143,0.1560693642) (0.9642857143,0.1560693642) (0.9642857143,0.1551724138) (0.9642857143,0.1551724138) (0.9642857143,0.1542857143) (0.9642857143,0.1542857143) (0.9642857143,0.1534090909) (0.9642857143,0.1534090909) (0.9642857143,0.1525423729) (0.9642857143,0.1525423729) (0.9642857143,0.1516853933) (0.9642857143,0.1516853933) (0.9642857143,0.1508379888) (0.9642857143,0.1508379888) (0.9642857143,0.15) (0.9642857143,0.15) (0.9642857143,0.1491712707) (0.9642857143,0.1491712707) (0.9642857143,0.1483516484) (0.9642857143,0.1483516484) (0.9642857143,0.1475409836) (0.9642857143,0.1475409836) (0.9642857143,0.1467391304) (0.9642857143,0.1467391304) (0.9642857143,0.1459459459) (0.9642857143,0.1459459459) (0.9642857143,0.1451612903) (0.9642857143,0.1451612903) (0.9642857143,0.1443850267) (0.9642857143,0.1443850267) (0.9642857143,0.1436170213) (0.9642857143,0.1436170213) (0.9642857143,0.1428571429) (0.9642857143,0.1428571429) (0.9642857143,0.1421052632) (0.9642857143,0.1421052632) (0.9642857143,0.1413612565) (0.9642857143,0.1413612565) (0.9642857143,0.140625) (0.9642857143,0.140625) (0.9642857143,0.1398963731) (0.9642857143,0.1398963731) (0.9642857143,0.1391752577) (0.9642857143,0.1391752577) (0.9642857143,0.1384615385) (0.9642857143,0.1384615385) (0.9642857143,0.137755102) (0.9642857143,0.137755102) (0.9642857143,0.1370558376) (0.9642857143,0.1370558376) (0.9642857143,0.1363636364) (0.9642857143,0.1363636364) (0.9642857143,0.135678392) (0.9642857143,0.135678392) (0.9642857143,0.135) (0.9642857143,0.135) (0.9642857143,0.1343283582) (0.9642857143,0.1343283582) (0.9642857143,0.1336633663) (0.9642857143,0.1336633663) (0.9642857143,0.1330049261) (0.9642857143,0.1330049261) (0.9642857143,0.1323529412) (0.9642857143,0.1323529412) (0.9642857143,0.1317073171) (0.9642857143,0.1317073171) (0.9642857143,0.1310679612) (0.9642857143,0.1310679612) (0.9642857143,0.1304347826) (0.9642857143,0.1304347826) (0.9642857143,0.1298076923) (0.9642857143,0.1298076923) (0.9642857143,0.1291866029) (0.9642857143,0.1291866029) (0.9642857143,0.1285714286) (0.9642857143,0.1285714286) (0.9642857143,0.1279620853) (0.9642857143,0.1279620853) (0.9642857143,0.1273584906) (0.9642857143,0.1273584906) (0.9642857143,0.1267605634) (0.9642857143,0.1267605634) (0.9642857143,0.1261682243) (0.9642857143,0.1261682243) (0.9642857143,0.1255813953) (0.9642857143,0.1255813953) (0.9642857143,0.125) (0.9642857143,0.125) (0.9642857143,0.1244239631) (0.9642857143,0.1244239631) (0.9642857143,0.123853211) (0.9642857143,0.123853211) (0.9642857143,0.1232876712) (0.9642857143,0.1232876712) (0.9642857143,0.1227272727) (0.9642857143,0.1227272727) (0.9642857143,0.1221719457) (0.9642857143,0.1221719457) (0.9642857143,0.1216216216) (0.9642857143,0.1216216216) (0.9642857143,0.1210762332) (0.9642857143,0.1210762332) (0.9642857143,0.1205357143) (0.9642857143,0.1205357143) (0.9642857143,0.12) (0.9642857143,0.12) (0.9642857143,0.1194690265) (0.9642857143,0.1194690265) (0.9642857143,0.1189427313) (0.9642857143,0.1189427313) (0.9642857143,0.1184210526) (0.9642857143,0.1184210526) (0.9642857143,0.1179039301) (0.9642857143,0.1179039301) (0.9642857143,0.1173913043) (0.9642857143,0.1173913043) (0.9642857143,0.1168831169) (0.9642857143,0.1168831169) (0.9642857143,0.1163793103) (0.9642857143,0.1163793103) (0.9642857143,0.1158798283) (0.9642857143,0.1158798283) (0.9642857143,0.1153846154) (0.9642857143,0.1153846154) (0.9642857143,0.114893617) (0.9642857143,0.114893617) (0.9642857143,0.1144067797) (0.9642857143,0.1144067797) (0.9642857143,0.1139240506) (0.9642857143,0.1139240506) (0.9642857143,0.1134453782) (0.9642857143,0.1134453782) (0.9642857143,0.1129707113) (0.9642857143,0.1129707113) (0.9642857143,0.1125) (0.9642857143,0.1125) (0.9642857143,0.112033195) (0.9642857143,0.112033195) (0.9642857143,0.1115702479) (0.9642857143,0.1115702479) (0.9642857143,0.1111111111) (0.9642857143,0.1111111111) (0.9642857143,0.1106557377) (0.9642857143,0.1106557377) (0.9642857143,0.1102040816) (0.9642857143,0.1102040816) (0.9642857143,0.1138211382) (1,0.1138211382) (1,0.1133603239) (1,0.1133603239) (1,0.1129032258) (1,0.1129032258) (1,0.1124497992) (1,0.1124497992) (1,0.112) (1,0.112) (1,0.1115537849) (1,0.1115537849) (1,0.1111111111) (1,0.1111111111) (1,0.1106719368) (1,0.1106719368) (1,0.1102362205) (1,0.1102362205) (1,0.1098039216) (1,0.1098039216) (1,0.109375) (1,0.109375) (1,0.1089494163) (1,0.1089494163) (1,0.1085271318) (1,0.1085271318) (1,0.1081081081) (1,0.1081081081) (1,0.1076923077) (1,0.1076923077) (1,0.1072796935) (1,0.1072796935) (1,0.106870229) (1,0.106870229) (1,0.1064638783) (1,0.1064638783) (1,0.1060606061) (1,0.1060606061) (1,0.1056603774) (1,0.1056603774) (1,0.1052631579) (1,0.1052631579) (1,0.1048689139) (1,0.1048689139) (1,0.1044776119) (1,0.1044776119) (1,0.1040892193) (1,0.1040892193) (1,0.1037037037) (1,0.1037037037) (1,0.1033210332) (1,0.1033210332) (1,0.1029411765) (1,0.1029411765) (1,0.1025641026) (1,0.1025641026) (1,0.102189781) (1,0.102189781) (1,0.1018181818) (1,0.1018181818) (1,0.1014492754) (1,0.1014492754) (1,0.1010830325) (1,0.1010830325) (1,0.1007194245) (1,0.1007194245) (1,0.1003584229) (1,0.1003584229) (1,0.1) (1,0.1) (1,0.09964412811) (1,0.09964412811) (1,0.09929078014) (1,0.09929078014) (1,0.09893992933) (1,0.09893992933) (1,0.0985915493) (1,0.0985915493) (1,0.09824561404) (1,0.09824561404) (1,0.0979020979) (1,0.0979020979) (1,0.09756097561) (1,0.09756097561) (1,0.09722222222) (1,0.09722222222) (1,0.09688581315) (1,0.09688581315) (1,0.09655172414) (1,0.09655172414) (1,0.09621993127) (1,0.09621993127) (1,0.09589041096) (1,0.09589041096) (1,0.09556313993) (1,0.09556313993) (1,0.09523809524) (1,0.09523809524) (1,0.09491525424) (1,0.09491525424) (1,0.09459459459) (1,0.09459459459) (1,0.09427609428) (1,0.09427609428) (1,0.09395973154) (1,0.09395973154) (1,0.09364548495) (1,0.09364548495) (1,0.09333333333) (1,0.09333333333) (1,0.09302325581) (1,0.09302325581) (1,0.09271523179) (1,0.09271523179) (1,0.09240924092) (1,0.09240924092) (1,0.09210526316) (1,0.09210526316)};
\addlegendentry{Pooled AP = 0.578}
\addplot[evalorange!70,dashed,line width=0.7pt] coordinates {(0,0.09210526316) (1,0.09210526316)};
\addlegendentry{ET prevalence = 0.092}
\addplot[forget plot,only marks,mark=*,mark size=1.8pt,evalorange] coordinates {(0.3571428571,0.7142857143)};
\node[anchor=west,font=\scriptsize,text=evalorange] at (axis cs:0.43214286,0.58428571) {$s>0$};
\draw[evalorange,thin,->] (axis cs:0.44214286,0.63428571) -- (axis cs:0.35714286,0.71428571);
\end{groupplot}
\end{tikzpicture}

\caption{외부 점수의 ROC와 ET 정밀도--재현율 곡선.
회색은 각 외부 폴드, 청록색은 참가자별 외부 점수의 통합 곡선이다.
주황색 점은 사전에 고정한 결정 규칙 \(s>0\)의 통합 작동점이다.
ROC 점선은 무정보 대각선, PR 점선은 전체 ET 비율
\(0.0921\)이다. PR 계단의 면적은 비보간 AP 정의와 일치한다.
다른 폴드 모델의 점수 척도가 다를 수 있어 통합 순위 지표는 폴드별
지표와 함께 해석한다.}\label{fig:evaluation-curves}
\end{figure}

통합 AUROC는 \(0.9000\), 통합 ET AP는 \(0.5776\)이다.
외부 폴드별 AUROC의 단순 평균은 \(0.9068\), 표본 표준편차는
\(0.0854\)이며, AP의 평균과 표본 표준편차는 각각
\(0.6684\), \(0.2631\)이다.
통합 지표는 서로 다른 모델 사이의 점수까지 비교하지만 폴드별 지표는
각 모델 내부의 순위만 비교하므로 두 집계 방식의 값이 같을 필요는 없다.
이 표준편차는 겹치는 학습 집합에서 얻은 기술적 산포이며 신뢰구간이 아니다.

순위 성능과 고정 임계값의 ET 검출률은 별개의 성질이다.
예를 들어 외부 폴드 4의 AUROC와 AP는 모두 \(1\)이지만,
임계값 \(0\)에서는 ET 6명 중 2명만 ET로 예측하였다.
PD보다 높은 점수를 받은 ET라도 점수 자체가 \(0\) 이하일 수 있기 때문이다.
따라서 높은 AUROC를 높은 ET 재현율 또는 확률 보정의 근거로 바꾸어 읽으면
안 된다. 이 결과를 본 뒤 외부 자료에서 임계값을 골라 같은 자료의 성능을
확증적 결과로 보고하는 것도 피해야 한다.

\clearpage
\subsection{외부 폴드의 인원과 정규화 선택}
\begin{table}[H]
\centering\small
\caption{외부 테스트 폴드의 클래스별 인원, ET 정답 수와
내부 검증에서 선택한 릿지 강도.}
\label{tab:observed-folds}
\begin{tabular}{@{}rrrrr@{}}
\toprule 외부 폴드 & PD 인원 & ET 인원 & ET 정답/전체 & 선택 \(\alpha\)\\ \midrule
1 & 56 & 5 & 2/5 & 1000\\
2 & 55 & 6 & 3/6 & 1000\\
3 & 55 & 6 & 2/6 & 1000\\
4 & 55 & 6 & 2/6 & 1000\\
5 & 55 & 5 & 1/5 & 1000\\

\bottomrule
\end{tabular}
\end{table}
다섯 외부 폴드 모두 탐색 격자의 상한인 \(\alpha=1000\)을 선택하였다.
이는 저장된 후보 안에서의 선택을 뜻하며 \(1000\)보다 큰 값이 최적인지는
이 실험으로 알 수 없다. 더 넓은 격자를 비교하려면 별도의 평가 계획과
중첩 선택 절차를 적용해야 한다.

\subsection{폴드별 지표의 산포}
\begin{figure}[H]
\centering
\begin{tikzpicture}
\begin{groupplot}[
 group style={group size=3 by 2,horizontal sep=0.9cm,vertical sep=1.2cm},
 width=4.55cm,height=4.3cm,xmin=0.6,xmax=5.4,ymin=-0.03,ymax=1.04,
 xtick={1,2,3,4,5},ytick={0,0.25,0.5,0.75,1},
 xlabel={Outer fold},axis lines=left,axis line style={evalgray},
 grid=major,grid style={black!8},
 tick label style={font=\scriptsize},label style={font=\footnotesize},
 title style={font=\footnotesize\bfseries},clip=true]
\nextgroupplot[title={(a) Balanced accuracy},xlabel={}]
\addplot[evalgray,densely dotted,line width=0.7pt] coordinates {(0.6,0.5) (5.4,0.5)};
\addplot[evalorange,dashed,line width=0.9pt] coordinates {(0.6,0.6713250518) (5.4,0.6713250518)};
\addplot[only marks,mark=*,mark size=2pt,evalnavy] coordinates {(1,0.6910714286) (2,0.75) (3,0.6575757576) (4,0.6666666667) (5,0.5818181818)};
\nextgroupplot[title={(b) ET recall},xlabel={}]
\addplot[evalgray,densely dotted,line width=0.7pt] coordinates {(0.6,0) (5.4,0)};
\addplot[evalorange,dashed,line width=0.9pt] coordinates {(0.6,0.3571428571) (5.4,0.3571428571)};
\addplot[only marks,mark=*,mark size=2pt,evalnavy] coordinates {(1,0.4) (2,0.5) (3,0.3333333333) (4,0.3333333333) (5,0.2)};
\nextgroupplot[title={(c) PD recall},xlabel={}]
\addplot[evalgray,densely dotted,line width=0.7pt] coordinates {(0.6,1) (5.4,1)};
\addplot[evalorange,dashed,line width=0.9pt] coordinates {(0.6,0.9855072464) (5.4,0.9855072464)};
\addplot[only marks,mark=*,mark size=2pt,evalnavy] coordinates {(1,0.9821428571) (2,1) (3,0.9818181818) (4,1) (5,0.9636363636)};
\nextgroupplot[title={(d) Accuracy}]
\addplot[evalgray,densely dotted,line width=0.7pt] coordinates {(0.6,0.9078947368) (5.4,0.9078947368)};
\addplot[evalorange,dashed,line width=0.9pt] coordinates {(0.6,0.9276315789) (5.4,0.9276315789)};
\addplot[only marks,mark=*,mark size=2pt,evalnavy] coordinates {(1,0.9344262295) (2,0.9508196721) (3,0.9180327869) (4,0.9344262295) (5,0.9)};
\nextgroupplot[title={(e) AUROC}]
\addplot[evalgray,densely dotted,line width=0.7pt] coordinates {(0.6,0.5) (5.4,0.5)};
\addplot[evalorange,dashed,line width=0.9pt] coordinates {(0.6,0.8999741201) (5.4,0.8999741201)};
\addplot[only marks,mark=*,mark size=2pt,evalnavy] coordinates {(1,0.9678571429) (2,0.8424242424) (3,0.9272727273) (4,1) (5,0.7963636364)};
\nextgroupplot[title={(f) ET average precision}]
\addplot[evalgray,densely dotted,line width=0.7pt] coordinates {(0.6,0.09210526316) (5.4,0.09210526316)};
\addplot[evalorange,dashed,line width=0.9pt] coordinates {(0.6,0.5776189524) (5.4,0.5776189524)};
\addplot[only marks,mark=*,mark size=2pt,evalnavy] coordinates {(1,0.7087301587) (2,0.6959259259) (3,0.6746031746) (4,1) (5,0.2627679567)};
\end{groupplot}
\end{tikzpicture}

\caption{다섯 외부 테스트 폴드의 지표. 점은 개별 폴드, 주황색 점선은
통합 외부 지표, 회색 점선은 전체 코호트의 항상 PD 기준선이다.
일반 정확도와 AP의 기준선은 폴드마다 조금 다른 클래스 비율을 보정한 값이
아니라 전체 코호트 기준값이다. 모든 패널의 세로축은 \(0\)--\(1\) 척도이다.
폴드를 잇는 추세선이나 신뢰구간은 표시하지 않는다.}
\label{fig:evaluation-folds}
\end{figure}

폴드별 균형 정확도는 \(0.5818\)--\(0.7500\)이고,
단순 평균은 \(0.6694\), 표본 표준편차는 \(0.0608\)이다.
폴드별 ET 재현율은 \(0.2000\)--\(0.5000\)이다.
ET가 각 폴드에 5명 또는 6명뿐이므로 한 사람의 예측 변화가 재현율을
20 또는 약 16.67 퍼센트포인트 바꾼다.
한 번의 5폴드 평가를 독립 실험 다섯 번으로 해석하지 않는다.

\paragraph{재현성.}
그림의 모든 수치 좌표를 이 원고에 직접 포함하였다.
지표 재계산, 입력 해시와 그림 재생성 절차는 부록~\ref{app:evaluation-rebuild}에 제시한다.
\clearpage
% END GENERATED OUTER EVALUATION

\section{실행 및 결과 파일}
\subsection{재현 환경과 실행}
프로젝트 루트에서 Python 3.12의 별도 환경을 만든다.
아래 잠금 파일은 검증 환경의 의존 버전을 기록한다.
\begin{lstlisting}
python3.12 -m venv .venv-sota
.venv-sota/bin/python -m pip install --upgrade pip
.venv-sota/bin/python -m pip install \
  -r sota_pd_et/requirements-lock.txt

.venv-sota/bin/python -m sota_pd_et.evaluate \
  --data-root data/pads/1.0.0 \
  --output artifacts/pads_mr_hydra \
  --outer-folds 5 --inner-folds 3 --repeats 1 \
  --n-jobs 4 --batch-size 32 --seed 20261017 \
  --fit-final
\end{lstlisting}
출력은 존재하지 않는 새 디렉터리를 지정한다.
예시 경로가 이미 있으면 다른 새 이름을 쓴다.
현재 구현은 중단된 실행을 자동 재개하지 않는다.
기본 계산은 CPU에서 수행하고 HYDRA에 PyTorch를 사용한다.
\code{batch-size}는 한 번에 변환할 참가자 수이며
확률적 경사 하강법의 학습 미니배치 크기가 아니다.
메모리가 부족하면 8 또는 16으로 낮출 수 있다.

\begin{table}[H]
\centering\small
\caption{주요 코드의 역할.}
\begin{tabularx}{\linewidth}{@{}lX@{}}
\toprule 파일 & 역할\\ \midrule
\code{data.py} & 진단 문자열 선택, 무결성 검사, 입력 행렬 구성.\\
\code{model.py} & 특징 추출ㆍ정규화, 가중 릿지와 예측 규칙.\\
\code{evaluate.py} & 내부 초모수 선택, 외부 평가, 실행 이력 저장.\\
\code{predict.py} & 저장 모델과 새 전처리 배열로 예측.\\
\bottomrule
\end{tabularx}
\end{table}

\Needspace{17\baselineskip}
\begin{longtable}{@{}p{0.37\linewidth}p{0.58\linewidth}@{}}
\caption{주요 출력 파일.}\label{tab:outputs}\\
\toprule 파일 & 의미\\ \midrule
\endfirsthead
\toprule 파일 & 의미\\ \midrule
\endhead
\code{summary.json} & 외부 평가 지표, 반복ㆍ폴드 요약, 항상 PD 기준선.\\
\code{oof_predictions.csv} & 참가자, 반복, 외부 폴드, 실제 라벨, 점수, 예측.\\
\code{fold_metrics.csv} & 외부 폴드별 지표와 선택된 \(\alpha\).\\
\code{splits.json} & 정확한 외부ㆍ내부 참가자 분할.\\
\code{cohort.csv} & 참가자 라벨과 감사용 표시.\\
\code{manifest.json} & 상태, 설정, 패키지 버전, 코드 및 자료 해시.\\
각 폴드의 \code{alpha_search.csv} & 내부 선택 지표. 외부 테스트 지표와 구분.\\
\code{final_model.joblib} & 요청 시 전체 자료로 다시 적합한 최종 모델.\\
\bottomrule
\end{longtable}

\subsection{최종 모델과 새 참가자 예측}
\code{--fit-final}을 지정하면 전체 304명에서 내부 3폴드 선택을 다시 수행하고,
선택한 \(\alpha\)로 전체 자료를 적합한다.
추가 특징 적합은 \(3+1=4\)번, 릿지 적합은 \(3\times7+1=22\)번이다.
이 최종 모델은 이후 예측을 위한 추정기이며,
이미 학습한 304명의 예측을 독립 테스트로 보고하면 안 된다.
새 입력은 동일한 전처리ㆍ채널 순서의 \((M,132,976)\) NumPy 배열이며 \(M\geq1\)이다.
\begin{lstlisting}
.venv-sota/bin/python -m sota_pd_et.predict \
  --model artifacts/pads_mr_hydra/final_model.joblib \
  --inputs new_participants.npy \
  --output artifacts/new_predictions.csv
\end{lstlisting}
원시 7열 기록을 그대로 넣는 형식은 아니다.
예측기는 저장된 주요 패키지 버전과 현재 환경의 일치 여부를 확인한다.

\section{문헌 근거와 해석의 한계}
\subsection{SOTA를 주장할 수 있는 범위}
MultiROCKET--HYDRA는 일반 시계열 분류 비교에서 경쟁력 있는 방법으로
보고되었다 \cite{bakeoff}. 이는 검토할 모델을 선택하는 근거이다.
그러나 그 자료 구성, 단변량ㆍ다변량 설정, 지표가 이 PD/ET 코호트와 같지는 않다.
2026년 현재 PD/ET 최고 성능이라고 하려면
동일한 참가자 분할ㆍ정보 범위ㆍ튜닝 예산으로 대조 모델을 비교해야 한다.

\subsection{점수는 확률이 아니며 계수는 인과적 설명이 아니다}
릿지 점수 \(s\)는 임의의 실수이며
식~\eqref{eq:ridge}는 \(s=P(y=1\mid X)\)를 보장하지 않는다.
예를 들어 \(s=0.8\)은 ET 확률 80\%라는 뜻이 아니다.
확률이 필요하면 학습 자료 내부의 별도 보정과 외부 보정 평가를 설계해야 한다.
현재 구현은 Brier 점수, 로그 손실, 보정 곡선을 보고하지 않는다.

특징은 여러 채널의 혼합과 비선형 요약이다.
큰 릿지 계수만으로 특정 손목ㆍ주파수ㆍ질병 기전이 인과적으로 중요하다고
결론 내릴 수 없다. 독립 기관, 다른 장비 또는 다른 진단 구성에서의
성능도 현재의 내부 코호트 평가만으로 알 수 없다.

\subsection{자료 선택과 검증 상태}
공식 전처리는 원 연구에서 이미 선택한 과제와 보정 규약을 포함한다.
참가자를 올바르게 분리하더라도 전처리 설계까지 완전히 미접촉인
전향적 실험이 되지는 않는다. 원시 시간 간격, 라벨의 모호성과
28명 ET 표본의 불확실성도 별도로 고려해야 한다.

코드 준비 단계에서 합성 자료의 12개 검사가 통과했고
실제 선택 코호트의 611개 원본 해시와 기본 분할을 확인하였다.
합성 통합 검사는 특징 변환, 중첩 평가, 저장ㆍ재로딩, 예측 진입점을 확인한다.
이 검증은 실제 PD/ET 진단 성능 측정과 별개이다.
실측 수치는 완료된 실행의 \code{summary.json}과 외부 예측을 대조하여
\ref{sec:observed-results}절에 별도로 제시하였다.
이 외부 교차검증은 독립 기관이나 전향적 진료 환경의 평가를 대신하지 않는다.

\clearpage
\appendix
\section{LSPV의 고정 버전 인덱스 규약}\label{app:lspv}
응답 \(v_0,\ldots,v_{m-1}\), 임계값 \(b\)를 고정한다.
현재 다변량 풀링 코드의 계산은 다음 순서와 같다.
\begin{enumerate}
\item 경계 인덱스 \(a\gets0\), 최댓값 \(L\gets0\)으로 초기화한다.
\item \(t=0,\ldots,m-1\)을 순서대로 확인한다.
 \(v_t<b\)이면 \(L\gets\max(L,t-a)\), 이어 \(a\gets t\)로 갱신한다.
 \(v_t\geq b\)이면 두 변수를 바꾸지 않는다.
\item 마지막으로 \(L\gets\max(L,m-1-a)\)로 갱신해 반환한다.
\end{enumerate}
PPV의 선택 조건 \(v_t>b\)와 위 경계 갱신 조건 \(v_t<b\)는 다르다.
동률과 끝점 때문에 단순한 연속 양수 개수와 차이가 생긴다.
예를 들어 모든 값이 \(b\)보다 크면 \(L=m-1\)이고,
\(m\geq2\)이며 모든 값이 \(b\)보다 작으면 \(L=1\)이다.
이는 별도의 이론적 지속시간을 정의하는 예가 아니라
고정 구현을 재현하는 데 필요한 규약이다.

\section{결과를 읽는 순서}
\begin{enumerate}
\item \code{manifest.json}에서 완료 상태, 버전, 설정, 시드를 확인한다.
\item 코호트의 진단 문자열과 클래스별 인원을 확인한다.
\item 반복마다 참가자가 외부 예측에 정확히 한 번씩 등장하는지 확인한다.
\item 균형 정확도, ET/PD 재현율, 혼동행렬을 먼저 읽는다.
\item 항상 PD 기준선과 비교하고 폴드별ㆍ통합 AUROC/AP를 함께 읽는다.
\item 내부 선택 지표와 외부 평가 지표, 전체 자료의 최종 모델을 구분한다.
\end{enumerate}

\section{평가 그림의 재생성}\label{app:evaluation-rebuild}
\code{sota_pd_et.plot_evaluation}은 완료된 실험의 출력만 읽어 지표를 다시 계산한다.
모델을 학습하거나 체크포인트를 불러오지 않는다.
외부 예측의 참가자별 유일성, 실제 라벨, 고정 결정 규칙,
외부ㆍ내부 분할 배정, 개별 폴드 파일 및 저장 지표의 일치를 확인한다.
AUROC는 ET/PD 점수 쌍의 비교로도 대조하고,
PR 계단의 면적이 비보간 AP와 같은지 확인한다.
실행 당시 기록된 평가 코드의 SHA-256과 현재 파일의 일치도 검사하였다.

\begin{lstlisting}
.venv-sota/bin/python -m sota_pd_et.plot_evaluation \
  --run artifacts/pads_mr_hydra --repeat 1 --update-article
\end{lstlisting}

그림 좌표, 재계산한 지표, 입력 파일 SHA-256과 검증 항목은
\code{output/figures/pads_mr_hydra/evaluation_data.json}에 저장한다.
같은 디렉터리의 세 독립형 그림 원고를 XeLaTeX로 컴파일하면 벡터 PDF를 얻는다.
이 문서에는 동일한 그림의 좌표를 직접 포함하므로
문서 컴파일에는 별도 그림 파일, 데이터 또는 모델 체크포인트가 필요하지 않다.
현재 재생성 문구는 이 304명ㆍ5폴드 실험에 맞춰져 있다.
코호트, 선택된 \(\alpha\) 또는 폴드 4의 예시가 달라지면
명령이 중단되어 관련 해석을 다시 작성하도록 한다.

\clearpage
\begin{thebibliography}{9}\small
\bibitem{pads}
PADS 연구진. \emph{PADS -- Parkinsons Disease Smartwatch Dataset},
version 1.0.0. PhysioNet, 2024.
\url{https://physionet.org/content/parkinsons-disease-smartwatch/1.0.0/}.
DOI: \href{https://doi.org/10.13026/m0w9-zx22}{10.13026/m0w9-zx22}.

\bibitem{multirocket}
C.~W.~Tan, A.~Dempster, C.~Bergmeir, G.~I.~Webb.
MultiRocket: Multiple pooling operators and transformations for fast and
effective time series classification.
\emph{Data Mining and Knowledge Discovery}, 2022.
\href{https://doi.org/10.1007/s10618-022-00844-1}{doi:10.1007/s10618-022-00844-1}.
\url{https://arxiv.org/abs/2102.00457}.

\bibitem{hydra}
A.~Dempster, D.~F.~Schmidt, G.~I.~Webb.
Hydra: competing convolutional kernels for fast and accurate time series classification.
\emph{Data Mining and Knowledge Discovery}, 37:1779--1805, 2023.
\href{https://doi.org/10.1007/s10618-023-00939-3}{doi:10.1007/s10618-023-00939-3}.

\bibitem{bakeoff}
M.~Middlehurst, P.~Sch\"afer, A.~Bagnall.
Bake off redux: a review and experimental evaluation of recent time series
classification algorithms. \emph{Data Mining and Knowledge Discovery}, 2024.
\href{https://doi.org/10.1007/s10618-024-01022-1}{doi:10.1007/s10618-024-01022-1}.
\url{https://arxiv.org/abs/2304.13029}.

\bibitem{aeonmr}
aeon 개발진. MultiROCKET 수치 구현, tag v1.6.0.
\url{https://github.com/aeon-toolkit/aeon/blob/v1.6.0/aeon/transformations/collection/convolution_based/_multirocket.py}.
이 문서의 수치 규약은 설치된 aeon 1.6.0의 소스에 따른다.

\bibitem{aeonhydra}
aeon 개발진. HYDRA 변환 및 희소 척도 구현, tag v1.6.0.
\url{https://github.com/aeon-toolkit/aeon/blob/v1.6.0/aeon/transformations/collection/convolution_based/_hydra.py}.
\url{https://github.com/aeon-toolkit/aeon/blob/v1.6.0/aeon/classification/convolution_based/_hydra.py}.

\bibitem{skridge}
scikit-learn 개발진. \emph{RidgeClassifier API documentation}.
\url{https://scikit-learn.org/stable/modules/generated/sklearn.linear_model.RidgeClassifier.html}.
참조 기준: scikit-learn 1.9.1.
\end{thebibliography}
\end{document}
